






















































































































































\section{Introduction}

A body $D\subseteq\mathbb{R}^n$ is said to be tiling the Euclidean space $\mathbb{R}^n$, if its translations by $\mathbb{Z}^n$
cover the entire space and have disjoint interiors. The foam problem asks for the least surface area a tiling body
$D$ can have. The problem had been considered by mathematicians already in the 19th century \cite{Thomson87},
and it also appears in chemistry, physics and engineering \cite{chemistry}. More recently, the problem had received
significant attention in the theoretical computer science community due to its strong relation with the parallel repetition problem~\cite{FKO,KROW,AK}.

The simplest example for a body that tiles the Euclidean space is the solid cube, $D = [0,1]^n$, which has surface
area $2n$. At first glance, one may expect the solid cube to be the best example there is, or more modestly
that any tiling body would need to have surface area $\Omega(n)$. The main results of \cite{KROW,AK}
show that this initial intuition is completely false, and that there are far more efficient tiling bodies whose surface area is $O(\sqrt{n})$.
This is surprising, since spheres --- which are the minimizers of surface area among all bodies with a given, fixed volume
(in this case volume $1$), have $\Theta(\sqrt{n})$ surface area and seem to be very far from forming a tiling of $\mathbb{R}^n$.
As we will shortly discuss, the existence of such surprising tiling body is intimately related to the existence of another
surprising object -- namely non-trivial strategies for $2$-prover-$1$-round games, repeated in parallel. The main goal of this
paper is to understand the symmetric variant of the foam problem, which is closely related to the symmetric variant of
parallel repetition.



















































































































\setcounter{thm}{3}\setcounter{subsection}{2}\subsection{Our results}
In this paper, our main object of study mainly are tilings of $\mathbb{R}^n$ using a \emph{symmetric body}.
\begin{definition}
  A set $D\subseteq\mathbb{R}^n$ is called symmetric if
  for any $\pi\in S_n$ and $x\in\mathbb{R}^n$, it holds that
  $x\in D$ if and only if $\pi(x)\in D$.
\end{definition}

The main question we consider, is what is the least surface area a symmetric tiling body can have.
Again, one has the trivial example of the solid cube $D = [0,1]^n$, but inspired by the non-symmetric
variant of the problem, one may expect there to be better examples. We first show that while this is
possible, the savings are much milder, and can be at most a multiplicative factor of $\sqrt{\log n}$.
\begin{thm}\label{thm:lb}
	Any symmetric tiling body $D$ of volume $1$
	with piecewise smooth surface has surface area at least $\Omega\left(\frac{n}{\sqrt{\log n}}\right)$.
\end{thm}

Besides the quantitative result itself, we believe the argument used in the proof of Theorem~\ref{thm:lb}
carries with it a lot of intuition regarding the additional challenge that the symmetric variants of the foam problem
and the parallel repetition posses, and we hope that this intuition will help us to develop better understanding
of symmetric parallel repetition in general. We remark that our proof actually shows a lower bound on the ``noise sensitivity''
parameter of the body, which is known to be smaller than the surface area of the body.

We complement Theorem~\ref{thm:lb} with a randomized construction showing that $O(\sqrt{\log n})$ savings
are indeed possible.
\begin{thm}\label{thm:ub}
	There exists a symmetric tiling body $D$ of volume $1$
	with piecewise smooth surface that has surface area $O\left(\frac{n}{\sqrt{\log n}}\right)$.
\end{thm}






































































































\setcounter{thm}{8}\section{Preliminaries}\label{sec:prelim}
\paragraph{Notations.}
We write $X\ll Y$ or $X = O(Y)$ to say that
there exists an absolute constant $C >0$ such that $X\leq C\cdot Y$,
and similarly write $X\gg Y$ or $X = \Omega(Y)$ to say that
there exists an absolute constant $c >0$ such that $X\geq c\cdot Y$.
We write $X\ee Y$ or $X = \Theta(Y)$ to say that $Y\ll X\ll Y$.


We denote random variables by boldface letters such as ${\bf x}$ and $\bm{\De}$.
We denote by $\mathcal{N}(\mu,\sigma^2)$ the distribution of a standard
Gaussian random variable with mean $\mu$ and variance $\sigma^2$, and by
$\mathcal{N}(\vec{\mu}, \Sigma)$ the distribution of a multi-dimensional
Gaussian random variable with means $\vec{\mu}$ and covariance matrix $\Sigma$.
\subsection{Needles}

\begin{definition}
  Let $\delta>0$, and let $a\in\mathbb{R}^n$. A random $\delta$-needle is a line $\ell_{a,{\bf u}} = \sett{a+t\cdot {\bf u}}{t\in[0,1]}$
  where the direction vector ${\bf u}$ is a chosen as a standard Gaussian $\mathcal{N}(0,\delta I_n)$.
\end{definition}

Given a tiling body $D$, a random $\delta$-needle from $D$ is a random $\delta$-needle $\ell_{{\bf a}, {\bf u}}$
where ${\bf a}\in D$ is chosen uniformly. Random needles are a useful tool to measure the surface area of a $D$, as shown
in the following two lemmas. First, given a tiling body $D$ and a needle $\ell_{a, u}$, we may think of the needle
as ``wrapping around'' around $D$,  i.e.\ its points are taken modulo $D$. We denote this ``wrapped around'' line by
$\tilde{\ell}_{a, u}$. We will use the following formula from \cite{Santalo}; the case $n=2$ is formula (8.10) therein,
and the extension to general $n$ is discussed in page 274.
\begin{lemma}\label{lem:santalo}
  There is a constant $C_n = \Theta(1)$, such that the following holds.
  Let $S$ be a piecewise smooth surface in a tiling body $D$ of volume $1$, and let $\delta>0$.
  Then
  \[
   \Expect{{\bf a}\in D, {\bf u}\sim\mathcal{N}(0,\delta I_n)}{\card{\tilde{\ell}_{{\bf a}, {\bf u}}\cap S}} = C_n \cdot\sqrt{\delta} \cdot {\sf area}(S).
  \]
\end{lemma}


\begin{lemma}\label{lem:area_bd}
  Let $D$ be a tiling body of volume $1$, and let $\delta>0$. Then
  \[
    \Prob{{\bf a}\in D, {\bf u}\sim\mathcal{N}(0,\delta I_n)}{\ell_{{\bf a},{\bf u}}\cap \partial D\neq \emptyset} \leq \Theta(\sqrt{\delta}) {\sf area}(\partial D).
  \]
\end{lemma}
\begin{proof}
  Set $S = \partial D$, and note that whenever $\ell_{a,\delta u}\cap \partial D\neq \emptyset$, we have that $\card{\tilde{\ell}_{a,\delta u}\cap S}\geq 1$.
  Hence by the previous lemma we get that
  \[
  \Prob{{\bf a}\in D, {\bf u}\sim\mathcal{N}(0,\delta I_n)}{\ell_{{\bf a},{\bf u}}\cap \partial D\neq \emptyset}
  \leq
  \Expect{{\bf a}\in D, {\bf u}\sim\mathcal{N}(0,\delta I_n)}{\card{\tilde{\ell}_{{\bf a},{\bf u}}\cap \partial D}}
  \leq  \Theta(\sqrt{\delta}) \cdot {\sf area}(\partial D).
  \qedhere
  \]
\end{proof}
We will use the above lemma to prove lower bounds on the surface area of a tiling body, by finding $\delta$ such that
the probability on the left hand side of Lemma~\ref{lem:area_bd} is at least $\Omega(1)$; this would imply that
${\sf area}(\partial D)\geq \Omega(1/\sqrt{\delta})$.

\subsection{Basic useful properties of tiling bodies}
\begin{lemma}\label{lem:diff}
  Let $D\subseteq \mathbb{R}^n$ be a symmetric body, such that for all $z\in\mathbb{Z}^n\setminus\set{0}$
  we have $D\cap (D+z)  = \emptyset$, and let $x\in D$. Then for every $1\leq i,j\leq n$,
  if $x_i - x_j \in \mathbb{Z}$, then $x_i = x_j$.
\end{lemma}
\begin{proof}
  Assume towards contradiction $x_i - x_j$ is a non-zero integer $k$, and let $S_{i,j}\in S_n$ be the permutation that
  maps $i$ to $j$, $j$ to $i$ and has any $r\neq i,j$ as a fixed point. Since $D$ is symmetric, we have
  that $S_{i,j}(x)\in D$. Also, we have
  \[
    x - S_{i,j}(x) = (x_i - x_j) (e_i - e_j) = k(e_i-e_j),
  \]
  where $e_i$ is the $i$th element in the standard basis. In other words, we get that $x = S_{i,j}(x) + z$ for non-zero $z\in \mathbb{Z}^n$,
  and therefore $x\in D + z$. This contradict the fact that $D$ and $D+z$ are disjoint.
\end{proof}

\begin{lemma}\label{lem:uniform}
  Let $D$ be a volume $1$ tiling body, and choose $a=(a_1,\ldots,a_n)\in D$ uniformly
  at random. Then the random variable $(a_1\hspace{-1ex}\pmod{1},\ldots,a_n\hspace{-1ex}\pmod{1})$ is uniform over $[0,1)^n$.
\end{lemma}
\begin{proof}
  Sample ${\bf x}\in[0,1)^n$, and take ${\bf a} = {\bf x} \pmod{D}$. Note that the distribution of
  ${\bf a}$ is uniform over $D$. Indeed, for that we note that the map $x\to x\pmod{D}$ is bijection from
  $[0,1)^n$ to $D$:
  otherwise, there were $x\neq x'$ in $[0,1)^n$ that are equal mod $D$, and therefore differ by non-zero lattice point
  (which is clearly impossible). Now as the distribution of ${\bf a}\pmod{1}$ is just ${\bf x}$, the claim follows.
  \end{proof}

\section{The lower bound: proof of Theorem~\ref{thm:lb}}\label{sec:lb}

In this section, we prove the lower bound on the surface area of a
symmetric tiling body $D$.
Throughout, we will have two parameters: $\sigma$, which is magnitude of each coordinates in the needle we consider
(which will be of order $\frac{\sqrt{\log n}}{n}$), and an auxiliary parameter $Z$ (which will be of order $\frac{n}{\log n}$).
Let $D$ be a symmetric tiling body containing $0$.
We denote by ${\bf a}$ a random point in $D$,
and by ${\bf u}$ a Gaussian vector $\mathcal{N}(0,\sigma^2 I_n)$. We will prove that
$\Prob{{\bf a}, {\bf u}}{\ell_{{\bf a},{\bf u}}\not\subseteq D} = \Omega(1)$,
which by Lemma~\ref{lem:area_bd} implies that ${\sf area}(\partial D)\geq \Omega(1/\sigma)$. As
$\sigma = \Theta(\sqrt{\log n}/n)$, this would establish Theorem \ref{thm:lb}.

\paragraph{Notations.}
For $x,y\in\mathbb{R}$, define
\[
d(x,y):= \min_{z\in \ZZ, z\neq 0} |(x+z)-y|\in [0,1].
\]
To gain some intuition for the definition of $d(x,y)$, suppose $x$ and $y$ are two entries of a point $a\in D$.
Clearly, if $d(x,y)$ is small, then $x,y$ nearly differ by an integer $z\neq 0$, and this says that the point $a$ is
somewhat close to the boundary of $D$ (in the sense that Lemma~\ref{lem:diff} could kick in if we move along a direction
that decreases this distance).


Our argument will indeed inspect $d(a_i,a_j)$ for all distinct $i,j\in[n]$ and the way they change along a random direction.
A key measure that we will keep track of is the energy of a point $a\in D$, defined by
\begin{equation*}
\Psi(a):= \sum_{i<j} e^{-Z \cdot d(a_i,a_j)}.
\end{equation*}
We show that for ${\bf a}\in_R D$ and ${\bf u}\sim \mathcal{N}(0,\sigma^2 I_n)$, if $\ell_{{\bf a}, {\bf u}}\subseteq D$ with
probability close to $1$, then the energy of ${\bf a}$ increases along the line $\ell_{{\bf a},{\bf u}}$ with high probability, and
in particular that $\Psi({\bf a}+{\bf u}) > \Psi({\bf a})$. We then argue that with high probability, this should be the case for
the point ${\bf a}$ as well as for ${\bf a} - {\bf u}$, hence $\Psi({\bf a}+{\bf u}) > \Psi({\bf a}-{\bf u})$ with high probability.
This event however can happen with probability at most $0.5$ by symmetry, hence completing the proof.


\subsection{Analyzing the energy along a random line}
By definition of $d(x,y)$, we either have $d(x,y) = (x+z-y)$ or $d(x,y) = -(x+z-y)$ for some $z\in\mathbb{Z}\setminus\set{0}$,
and this sign determines whether $x,y$ need to move in different directions or the same direction for $d(x,y)$ to get smaller.
To capture this, we denote
\[
\ga(x,y):=\left\{ \begin{array}{ll}
+1 & \text{ if $d(x,y)=x+z-y$ for some $z\in\ZZ$, $z\neq 0$},\\
-1 & \text{ otherwise}.	
\end{array} \right.
\]
Next, we discuss the {\em energy} of a configuration, which is the key concept used in the proof.
Let $Z$ be a parameter to be chosen later (of the order $n/\log n$).
As stated earlier, our goal is to analyze the behaviour of $\Psi(a)$ along a random $\sigma^2$-needle from $a$
in direction $u$. Towards this end, note that we expect (at least if $u_i, u_j$ are small) that $d(a_i+u_i, a_j+u_j) = d(a_i,a_j) + \ga(a_i,a_j) (u_j-u_j)$,
hence expect $\Psi(a+u)$ to be close to
\begin{equation*}
\Psi(a,u):=  \sum_{i<j} e^{-Z \cdot (d(a_i,a_j)+\ga(a_i,a_j)\cdot (u_i-u_j))}.
\end{equation*}
Indeed, this is the content of the following claim.

\begin{claim}
	\label{cl:1}
	Suppose $|u_i|\leq 1/20$ for all $i$, and $a+[0,1]\cdot u \subset D$, then
	$$
	|\Psi(a+u)-\Psi(a,u)|\leq n^2\cdot e^{-Z/4}.
	$$
\end{claim}

\begin{proof}
   We consider the contribution of each pair $(i,j)$ to $\Psi(a+u)$ and $\Psi(a,u)$ separately.
   Without loss of generality we may only consider pairs $i,j$ that $\ga(a_i,a_j)=1$,
   and thus $d(a_i,a_j)= a_i - a_j + z$ for some $z\in \ZZ$, $z\neq 0$.
   Let
   \[
   d= a_i - a_j + z + (u_i-u_j) = (a_i+u_i)- (a_j+u_j)+z.
   \]
   First, we argue that $d\geq 0$. Otherwise, since $a_i-a_j + z \geq 0$ it follows by continuity that there is
   $\lambda\in[0,1]$ such that $a_i - a_j + z + \lambda(u_i-u_j) = 0$, and hence the point $a+\lambda u$ has entries
   that differ by an integer $z\neq 0$, and this contradicts Lemma \ref{lem:diff} (as $a+\lambda u \in D$).
   We now consider two cases:
   \begin{itemize}
     \item {\bf Case 1}: $d \in [0,0.5]$. In this case, we have $d(a_i+u_i, a_j+u_j) = d$, and thus the contribution of the pair $(i,j)$ to both sums is the same ($e^{-Z\cdot d}$).
     \item {\bf Case 2}: $d > 0.5$. Since $\card{u_i-u_j}\leq 0.1$, it follows that $d(a_i,a_j) = d -(u_i-u_j) > 0.4$, which implies $d(a_i+u_i,a_j+u_j) > 0.3$. Therefore, the
     contribution to $\Psi(a,u)$ from $i,j$ is at most $e^{-0.4\cdot Z}$ and to $\Psi(a+u)$ is at most $e^{-0.3 \cdot Z}$, and in particular $(i,j)$ contributes (in absolute value)
     at most $e^{-Z/4}$ to the difference between the sums.
   \end{itemize}
   Taking a sum over all pairs $(i,j)$ concludes the proof.
\end{proof}


\subsection{Analyzing the expectation and variance of $\Psi(a,{\bf u})$}
Next,
we consider $\Psi(a,{\bf u})$ as a random variable over the choice of ${\bf u}$
and compute its expectation and variance. In both computations we will use
the well-known fact that $\E[e^{-Z\cdot N(0, c^2)}]=e^{Z^2 c^2/2}$ for all $c> 0$.

\begin{claim}
\label{cl:2}
For every $a\in\mathbb{R}^n$ we have $\Expect{{\bf u}\sim \mathcal{N}(0,\sigma^2 I_n)}{\Psi(a,{\bf u})} =\Psi(a)\cdot e^{(Z \cdot \si)^2}$.
\end{claim}

\begin{proof}
    By linearity of expectation we have that
	\[
	\Expect{{\bf u}\sim \mathcal{N}(0,\sigma^2I_n)}{\Psi(a,{\bf u})}
= \sum_{i<j} e^{-Z \cdot d(a_i,a_j)}\cdot \Expect{{\bf u}\sim \mathcal{N}(0,\sigma^2I_n)}{e^{-Z\cdot \ga(a_i,a_j)\cdot ({\bf u}_i-{\bf u}_j)}}.
	\]
    Note that the above expectation does not depend on $i,j$: for every $i,j$ the distribution of
    ${\bf u}_i - {\bf u}_j$ is $N(0,\si^2)-N(0,\si^2)\sim N(0,2\si^2)$, so it is symmetric around $0$ and thus
    the sign $\ga(a_i,a_j)$ does not affect the expectation. Hence we have
    \[
    \Expect{{\bf u}\sim \mathcal{N}(0,\sigma^2I_n)}{\Psi(a,{\bf u})}   = \Psi(a) \cdot \E [e^{Z\cdot N(0,2\sigma^2)}]
	=\Psi(a)\cdot e^{Z^2\si^2}.\qedhere
    \]
\end{proof}


Next, we turn our attention into upper bounding the variance of $\Psi(a,{\bf u})$, and
for that we first define the notion of {\em good} points $a\in D$ and prove two preliminary
claims.
We say a point $a$ is {\em good} if any interval of length $(10\log n)/n$
on the torus contains at least $\log n$ and at most $100 \log n$ coordinates from $a\hspace{-1ex}\pmod{1}$.
Note by Lemma \ref{lem:uniform}, if ${\bf a}$ is chosen randomly from $D$ then ${\bf a}\pmod{1}$ is uniform
over $[0,1)^n$ and by Chernoff bound is easily shown to be good with probability $>0.999$.

We first show that good points have high energy.
\begin{claim}\label{cl:6}
There exists $c_2>0$, such that for
$Z = 0.1 \frac{\log n}{n}$, if
$a$ is good then $\Psi(a)> c_2 \log^2 n$.
\end{claim}

\begin{proof}
 Partition the torus $[0,1)$ into $m = n/(10 \log n)$ disjoint intervals of length $1/m = (10 \log n)/n$ each.
 We say that $I_i$ is {\em unanimous}, if there is $b_i\in \RR$ (called anchor) such that
 (1) $b_i\hspace{-1ex}\pmod{1}$ is the middle of $I_i$, and (2) for the majority of points $a_j\in I_i$, $|a_j-b_i|< 1/m$.

 We consider two cases:

\noindent {\bf Case 1: There is an interval $I_i$ that is not unanimous.} Note that there are at least $\log n$ coordinates $j$ of $a$ such that $a_j\in I_i$.
Let $j^{\star}$ be such coordinate, and write $a_{j^{\star}} = z_{j^{\star}} + \{a_{j^{\star}}\}$ where $z_{j^{\star}}\in\mathbb{Z}$ and
$\{a_{j^{\star}}\}$ is the fractional part of $a_{j^{\star}}$. Consider
$b = z_{j^{\star}} + m_i$ where $m_i$ is the middle of $I_i$. Then since $I_i$ is not unanimous, $b$ is not an anchor of it and so there are at least $\half \log n$
coordinates of $a$, say $(a_k)_{k\in K_{i,j^{\star}}}$ that mod $1$ are in $I_i$, and $\card{a_k - b}\geq 1/m$. Writing $a_k = z_k + \{a_k\}$, we observe that $z_k\neq z_{j^{\star}}$,
since otherwise $\card{a_k - b}= \card{\{a_k\} - m_i} \leq 1/(2m)$. Hence the difference $a_k-a_{j^{\star}}$
is $10\log n/n$ close to an integer $z_k-z_{j^{\star}}\neq 0$, and so $d(a_k,a_{j^{\star}})\leq 10\log n/n$, and the contribution of $\Psi(a)$ is
at least $e^{-1}$. Summing we get
\[
\Psi(a)\geq
\half
\sum\limits_{j^{\star}: a_{j^{\star}} \in I_i}
\sum\limits_{k\in K_{i,j^{\star}}}{e^{-Z d(a_k,a_{j^{\star}})}}
\geq
\half
\sum\limits_{j: a_j \in I_i}
e^{-1} \card{K_{i,j^{\star}}}
\geq \frac{1}{4e} \log^2 n.
\]

\noindent {\bf Case 2: All intervals are unanimous.} Let $b_i$ be an anchor of $I_i$. Note that since the fractional part of two adjacent anchors,
i.e.\ of $b_i,b_{i+1}$, are $1/m$ apart, we have that either $\card{b_i - b_{i+1}}\leq 1/m$ or $\card{b_i-b_{i+1}}\geq 1-1/m$. We claim there exists $i$ for
which the latter condition holds. To see this, assume that for all $i=1,\ldots,m-1$ we have that the first condition holds. Then we have
$b_i = z + i \frac{10\log n}{n}$ for some $z\in \mathbb{Z}$ for all $i=1,\ldots,m$, and hence $\card{b_m - b_1}\geq 1-1/m$ (and the condition holds for $i=m$).

Thus, we fix $i$ such that $\card{b_i-b_{i+1}}\geq 1-1/m$, and thus $b_i - b_{i+1} = z + \alpha$ for $z\neq 0$ and $\card{\alpha}\leq 1/m$.
Let $K_i$ be the coordinates $j$ of $a$ such that $\card{a_j - b_i}\leq 1/m$ for $j\in K_i$
and similarly define $K_{i+1}$.
We have that $a_r - a_j = z + \alpha + (a_r - b_{i+1}) + (a_j - b_i)$, hence $a_r - a_j = z + \beta$ for $\card{\beta}\leq 3/m$ for all $r\in K_{i+1}$, $j\in K_i$.
Thus $d(a_r,a_j)\leq 3/m$, and we get
\[
\Psi(a)\geq \card{K_i}\card{K_{i+1}} e^{-Z\cdot 3/m}\geq \frac{1}{4} e^{-3} \log^2 n\qedhere
\]
\end{proof}


Let $C_i = \sum_{j\neq i} e^{-Z\cdot d(a_i,a_j)}$
be the contribution of $a_i$ to $\Psi(a)$. Note that $\Psi(a)= \frac{1}{2}\sum_i C_i$.

\begin{claim}\label{cl:7}
There exists $c_3>0$, such that if $a$ is good, then for all $i$ we have
$C_i < c_3 \Psi(a)/\log n$.
\end{claim}
\begin{proof}
	Note that $d(a_i,a_j)\ge | \{a_i\} - \{a_j\}|$. Since any interval of length $10\log n/n$ on the torus contains at most $100\log n$ points of $a$,
    we have that the number of $j$'s such that $|\{a_i\} - \{a_j\}|$ is between $10 \log n/n \cdot k$ and $10 \log n/n \cdot (k+1)$ is at most $200\log n$
    (for all $k$). Therefore,
	\[
	C_i < 200 \log n\cdot \sum_{k=0}^{\infty} e^{-Z \cdot k \cdot (10 \log n )/n} = 200 \log n \cdot \sum_{k=0}^{\infty}e^{-k}
    \leq 400 \log n.
	\]
Using Claim~\ref{cl:6}, we may bound $\log n\leq \frac{1}{c_2}\frac{\Psi(a)}{\log n}$, finishing the proof.
\end{proof}


We are now ready to bound the variance of $\Psi(a,{\bf u} )$.
\begin{claim}\label{cl:3}
There exists $c_1>0$ such that the following holds.
Let $Z= n/10\log n$, let $a\in \mathbb{R}^n$ be good and let ${\bf u} \sim \mathcal{N}(0,\sigma^2 I_n)$.
Then
\[
	{\sf var}_{\bf u} [\Psi(a,{\bf u} )] \leq \frac{c_1}{\log n} \cdot( e^{4(Z \cdot \si)^2}
	-e^{2(Z \cdot \si)^2})\cdot \Psi(a)^2.
\]
\end{claim}

\begin{proof}
Using Claim \ref{cl:2} to compute the expectation of $\Psi(a,{\bf u})$, we have by definition that
\begin{align*}
&{\sf var}_{\bf u}  ( \Psi(a,{\bf u} ))
= \E_{\bf u}  \left[\left(\sum_{i<j} e^{-Z\cdot d(a_i,a_j)}\cdot (e^{Z\cdot\gamma(a_i,a_j)\cdot({\bf u} _i-{\bf u} _j)}-e^{(Z\cdot \si)^2}) \right)^2\right]\\
& =\sum_{i<j} e^{-2Z\cdot d(a_i,a_j)} \cdot\Expect{{\bf u}}{\left(e^{Z\cdot\gamma(a_i,a_j)\cdot({\bf u}_i-{\bf u}_j)}-e^{(Z\cdot \si)^2}\right)^2}\\
& + \sum_{\substack{(i,j,k) \\ \text{distinct}}}e^{-Z\cdot (d(a_i,a_j)+d(a_i,a_k))}\cdot
\Expect{{\bf u}}{(e^{Z\cdot\gamma(a_i,a_j)\cdot({\bf u}_i-{\bf u}_j)}-e^{(Z\cdot \si)^2})(e^{Z\cdot\gamma(a_i,a_k)\cdot({\bf u}_i-{\bf u}_k)}-e^{(Z\cdot \si)^2})}.
\end{align*}
Here, we used that fact that if $i,j,k,r$ are all distinct then $e^{Z\cdot\gamma(a_i,a_j)\cdot({\bf u}_i-{\bf u}_j)}$, $e^{Z\cdot\gamma(a_k,a_r)\cdot({\bf u}_k-{\bf u}_r)}$
are independent with expectation $e^{(Z\cdot \sigma)^2}$, hence the contribution of these terms is $0$. Computing, we see that
\[
\Expect{{\bf u}}{\left(e^{Z\cdot\gamma(a_i,a_j)\cdot({\bf u}_i-{\bf u}_j)}-e^{(Z\cdot \si)^2}\right)^2}
=\E\left[e^{Z\cdot N(0,8\sigma^2)}\right]-e^{2(Z\cdot \si)^2}
=e^{4(Z\cdot \si)^2}-e^{2(Z\cdot \si)^2},
\]
and
\begin{align*}
&\Expect{{\bf u}}{(e^{Z\cdot\gamma(a_i,a_j)\cdot({\bf u}_i-{\bf u}_j)}-e^{(Z\cdot \si)^2})
 (e^{Z\cdot\gamma(a_i,a_k)\cdot({\bf u}_i-{\bf u}_k)}-e^{(Z\cdot \si)^2})}\\
 &=
 \E\left[e^{(\gamma(a_i,a_j) + \gamma(a_i,a_k))Z\cdot N(0,\sigma^2)}\right]
 \E\left[e^{Z\cdot N(0,2\sigma^2)}\right]
 -e^{2(Z\cdot \si)^2}\\
 &\leq
 \E\left[e^{2Z\cdot N(0,\sigma^2)}\right]
 \E\left[e^{Z\cdot N(0,2\sigma^2)}\right]
 -e^{2(Z\cdot \si)^2}\\
 &= e^{3(Z\cdot \sigma)^2}-e^{2(Z\cdot \si)^2}.
\end{align*}
Thus, we get that
\begin{align*}
  {\sf var}_{\bf u} ( \Psi(a,{\bf u}))
  &\leq
 \sum_{i<j} e^{-2Z\cdot d(a_i,a_j)} (e^{4 (Z\cdot \si)^2}-e^{2 (Z\cdot \si)^2}) +
 \sum_{\text{$(i,j,k)$ distinct}} e^{-Z\cdot (d(a_i,a_j)+d(a_i,a_k))}(e^{3(Z\cdot \sigma)^2}-e^{2(Z\cdot \si)^2}) \\
 & \leq
 (e^{4 (Z\cdot \si)^2}-e^{2 (Z\cdot \si)^2})
 \sum_{i}
 \left(\sum_{j\neq i} e^{-2Z\cdot d(a_i,a_j)} + \sum_{j,k\neq i}
e^{-Z\cdot (d(a_i,a_j)+d(a_i,a_k))}\right)\\
&=(e^{4 (Z\cdot \si)^2}-e^{2 (Z\cdot \si)^2})\sum\limits_{i}\left(\sum_{j\neq i} e^{-2Z\cdot d(a_i,a_j)}\right)^2\\
&=  (e^{4 (Z\cdot \si)^2}-e^{2 (Z\cdot \si)^2}) \cdot \sum_i C_i^2.
\end{align*}
Therefore, using Claim~\ref{cl:7} we conclude that
\begin{multline*}
{\sf var}_{\bf u} ( \Psi(a,{\bf u}))
\leq
( e^{4 (Z\cdot \si)^2}-e^{2 (Z\cdot \si)^2}) \frac{c_3\Psi(a)}{\log n}\cdot \sum_i C_i =
\frac{2 c_3}{\log n}\cdot
( e^{4 (Z\cdot \si)^2}-e^{2 (Z\cdot \si)^2}) \cdot \Psi(a)^2.
\end{multline*}
Setting $c_1:=2c_3$ completes the proof.
\end{proof}

Putting the last two claims together, we have:

\begin{claim}
	\label{cl:4}
Let $\sigma = 10^4\sqrt{c_1} \frac{\sqrt{\log n}}{n}$ and let $a\in\mathbb{R}^n$ be good.
Then
\[
\Pr_{{\bf u}}[\Psi(a,{\bf u})>\Psi(a)+\frac{(Z\si)^4}{2}\Psi(a)]\geq 0.96.
\]
\end{claim}
\begin{proof}
	We upper bound the probability of the complement event. Using Claim \ref{cl:2} (and $e^{t}\geq 1+t+t^2/2$ for $t\geq 0$), we get
	\[
	\E_{\bf u}[\Psi(a,{\bf u})]\geq \Psi(a)\cdot \left(1+(Z \si)^2+\frac{(Z\si)^4}{2}\right).
	\]
	Hence
	\[
	\Pr_{{\bf u}}\left[\Psi(a,{\bf u})\leq \Psi(a)+\frac{(Z\si)^4}{2}\Psi(a)\right] \leq \Pr_{{\bf u}}\left[\card{\Psi(a,{\bf u})- \E_{{\bf u'}}[\Psi(a,{\bf u'})]}\geq
\Psi(a)\cdot  (Z \si)^2\right].
	\]

    We want to upper bound the probability of the last event using Chebyshev's inequality.
    Since $a$ is good, the conclusion of Claim~\ref{cl:3} holds.
    Since $Z \si =o(1)$, for large enough $n$ we get
	\[
	{\sf var}_{\bf u} [\Psi(a,{\bf u})]
    \leq
    \frac{c_1}{\log n} \cdot( e^{4(Z \cdot \si)^2}
	-e^{2(Z \cdot \si)^2})\cdot \Psi(a)^2
    \leq \frac{c_1}{\log n} \cdot \Psi(a)^2\cdot 8(Z \si)^2.
	\]
    Therefore, applying Chebyshev's inequality we see the probability in question is at most
    \[
	\frac{{\sf var}_{\bf u} [\Psi(a,{\bf u})]}{\Psi(a)^2 \cdot  (Z \si)^4} \leq
	\frac{c_1 \cdot \Psi(a)^2 \cdot 8 (Z \si)^2}{(\log n)\cdot \Psi(a)^2 \cdot  (Z \si)^4 }
    = \frac{8c_1}{(\log n)\cdot (Z\sigma)^2 } = \frac{4 c_1}{10^2 c_1} = 0.04.\qedhere
    \]
\end{proof}


\subsection{Finishing the argument}
For each $u$, denote $\eps_u = \Prob{{\bf a}\in D}{\ell_{{\bf a},u}\not\subseteq D}$,
and denote $\eps = \Expect{{\bf u}\sim\mathcal{N}(0,\sigma^2 I_n)}{\eps_{\bf u}} = \Prob{{\bf a}, {\bf u}}{\ell_{{\bf a}, {\bf u}}\not\subseteq D}$.
\begin{claim}\label{cl:tv5}
For each $u$, $\mathcal{D}_{TV}[{\bf a}; {\bf a}-u]\leq \eps_u + \eps_{-u}$.
\end{claim}
\begin{proof}
  Let $K$ be a Borel set.
  Note that it is enough to show that
  (1) if $K\subseteq D$ then $0\leq \Prob{{\bf a}\in D}{{\bf a}\in K} - \Prob{{\bf a}\in D}{{\bf a}-u\in K}\leq \eps_u$,
  and (2) if $K\subseteq\bar{D}$, then $-\eps_{-u}\leq \Prob{{\bf a}\in D}{{\bf a}\in K} - \Prob{{\bf a}\in D}{{\bf a}-u\in K}\leq 0$.
  Indeed, given both (1) and (2), the triangle inequality implies for any Borel set $K\subseteq\mathbb{R}^n$,
  \begin{align*}
  &\card{\Prob{{\bf a}\in D}{{\bf a}\in K} - \Prob{{\bf a}\in D}{{\bf a}-u\in K}}\\
  &\leq
  \card{
  \Prob{{\bf a}\in D}{{\bf a}\in K\cap D}-\Prob{{\bf a}\in D}{{\bf a}-u\in K\cap D}
  +\Prob{{\bf a}\in D}{{\bf a}\in K\setminus D}-\Prob{{\bf a}\in D}{{\bf a}-u\in K\setminus D}}
  \leq \eps_u+\eps_{-u}.
  \end{align*}

  To prove (1), note that
  $\Prob{{\bf a}\in D}{{\bf a}\in K}=\mu(K)$
  and
  \[
  \Prob{{\bf a}\in D}{{\bf a}-u\in K} = \Prob{{\bf a}\in D}{{\bf a}\in K+u} = \mu((K+u)\cap D).
  \]
  This is at most $\mu(K+u) = \mu(K)$ (hence the expression in (1) is non-negative)
  and at least $\geq \mu(K+u) - \mu((K+u)\setminus D)= \mu(K) - \mu(K\setminus (D-u))$.
  Therefore
  \[
  0\leq \Prob{{\bf a}\in D}{{\bf a}\in K} - \Prob{{\bf a}\in D}{{\bf a}-u\in K}
  \leq \mu(K\setminus (D-u))
  \leq \mu(D\setminus (D-u))
  = \Prob{{\bf a}\in D}{{\bf a}+u\not\in D}
  \leq \eps_u.
  \]

  To prove (2), note that $\Prob{{\bf a}\in D}{{\bf a}\in K} = 0$ (hence the expression in (2) is non-positive)
  and
  \[
  \Prob{{\bf a}\in D}{{\bf a}-u\in K} \leq \Prob{{\bf a}\in D}{{\bf a}-u\not\in D} \leq \eps_{-u}.\qedhere
  \]
\end{proof}

\begin{claim}
	\label{cl:5}
	$\eps\geq 0.1$.
\end{claim}
\begin{proof}
	Let $E_1$ be the event that ${\bf a} + {\bf u}[0,1]\subseteq D$, let
    $E_2$ be the event that $\Psi({\bf a})\leq 1$, let $E_3$ be the event
	that $|{\bf u}_i|>1/20$ for some $i$ and let $E_4$ be the event that
	$\Psi({\bf a},{\bf u})>\Psi({\bf a})+\frac{(Z\si)^4}{2}\Psi({\bf a})$. Finally, let $E_5$ be the event that
	$\Psi({\bf a}+{\bf u})>\Psi({\bf a})$ and denote $E({\bf a},{\bf u}) = E_1\cap (\neg E_2)\cap (\neg E_3)\cap E_4$.
    Note that if the event $E$ holds for $a,u$, then $E_5$ also holds, since by Claim~\ref{cl:1}:
	\[
	\Psi(a+u)\geq \Psi(a,u)-n^2\cdot e^{-Z/4}> \Psi(a)+\frac{(Z\si)^4}{2}\Psi(a)-n^2\cdot e^{-Z/4}\geq \Psi(a).
	\]



    By Claim \ref{cl:6} the probability of $E_2$ is at most the probability ${\bf a}$ is bad, hence it is at most $0.005$.
    By definition, the probability of $E_1$ is $1-\eps$.
    By the union bound and Chernoff inequality, the probability of $E_3$ is $o(1)$.
	Thus, by Claim~\ref{cl:4} we have
	\begin{equation}
	\Pr_{{\bf a},{\bf u}}[E({\bf a},{\bf u})]\geq 0.96 -\eps- 0.005 - o(1)\geq 0.95-\eps.
	\end{equation}
    Fix $u$. Using Claim~\ref{cl:tv5} we get that
    \[
    \Prob{{\bf a}}{E({\bf a}-u,u)} \geq
    \Prob{{\bf a}}{E({\bf a},u)} - \mathcal{D}_{TV}[{\bf a}; {\bf a}-u]
    \geq \Prob{{\bf a}}{E({\bf a},u)} - \eps_u-\eps_{-u}.
    \]
    By the union bound, we now conclude that
    \[
    \Prob{{\bf a}}{E({\bf a}-u,u)\cap E({\bf a},u)}\geq
    1 -  \Prob{{\bf a}}{\overline{E({\bf a}-u,u)}}
    -  \Prob{{\bf a}}{\overline{E({\bf a},u)}}
    \geq 2 \Prob{{\bf a}}{E({\bf a},u)} -1-\eps_u-\eps_{-u}.
    \]
    Taking expectation over ${\bf u}$, we get that
    \[
      \Prob{{\bf a},{\bf u}}{E({\bf a}-{\bf u},{\bf u})\cap E({\bf a},{\bf u})}\geq
     2\Prob{{\bf a},{\bf u}}{E({\bf a},{\bf u})} -1-2\Expect{{\bf u}}{\eps_{\bf u}}
     \geq 0.9 - 4\eps.
    \]
    Next, when both $E(a-u,u)$ and $E(a,u)$ hold, we have by the previous observation that $E_5$ holds for both pairs $(a-u,u)$ and $(a,u)$,
    and so $\Psi(a+u) > \Psi (a) = \Psi((a-u) + u) > \Psi(a-u)$. Thus, we get that
    $\Prob{{\bf a},{\bf u}}{\Psi({\bf a}+{\bf u}) > \Psi({\bf a}-{\bf u})}\geq 0.9 - 4\eps$. On the other hand, the probability on the left hand side is at most $0.5$;
    this follows as
    $\Prob{{\bf a},{\bf u}}{\Psi({\bf a}+{\bf u}) > \Psi({\bf a}-{\bf u})}
    = \Prob{{\bf a},{\bf u}}{\Psi({\bf a}-{\bf u}) > \Psi({\bf a}+{\bf u})}$ (since the distributions of ${\bf u}$ and $-{\bf u}$ are identical) and their sum is at most $1$.
    Combining the two inequalities we get that $\eps \geq 0.1$.
\end{proof}


\setcounter{thm}{21}\setcounter{section}{3}\section{The upper bound: proof of Theorem~\ref{thm:ub}}\label{sec:ub}
In this section we prove a matching upper bound on the surface area of a symmetric foam by giving a
(probabilistic) construction of a symmetric tiling body $D$ of surface area $O(n/\sqrt{\log n})$.
The main technical result proved in this section, Lemma~\ref{lem:main_ub}, establishes a weaker statement,
and in Section~\ref{sec:ns_sa_reduce} we show how to deduce Theorem~\ref{thm:ub} from it.

\subsection{Reduction to constructing a rounding scheme}\label{sec:mapping_to_tiling}
Suppose $S$ is function mapping (multi-)sets of $n$ points from $\RR/\ZZ$, to $\RR/\ZZ$.
We further assume that for all (multi-)sets $A$, it holds that $S(A)\not\in\set{0}\cup A$.


Given such $S$, we may extend it to $\mathbb{R}^n$ by
$S(x_1,\ldots,x_n):=S(\{\{x_1\},\ldots,\{x_n\}\})$, where $\{x_i\}$ is the fractional part of $x$.
We can construct a rounding scheme
$R\colon \RR^n\to \ZZ^n$ using $S$ as follows.
\begin{itemize}
	\item
	On input $x=(x_1,\ldots,x_n)$, denote $z=S(x)$ and view
	$z$ as a number in $[0,1)$.
	\item For each $i\in [n]$:
\begin{itemize}
    \item
	if $\{x_i\}\in [0,z)$, set $R(x)_i = \lfloor x_i \rfloor$,
    \item otherwise, $\{x_i\}\in (z,1)$, and set $R(x)_i = \lceil x_i \rceil$.
\end{itemize}
\end{itemize}
First, $R$ is well-defined since $z\notin \{0,\{x_1\},\ldots,\{x_n\}\}$.
Next, note that for any $t\in\ZZ^n$ it holds that $R(x+t)=R(x)+t$, thus $R$ induces
that the body $D = \sett{x}{R(x) = 0}$ is tiling with respect to the lattice $\ZZ^n$.
Last, we note that since for any $\pi\in S_n$ we have that $S(\pi(x))=S(x)$, we also have
that $R(\pi(x))=\pi(R(x))$, and hence $D$ is symmetric.

In our proof we will define a distribution over mappings $S$, and we will want to study
the noise sensitivity of the resulting body $D$ using properties of the mappings $S$.
The following claim gives useful conditions to study noise sensitivity in terms of
mapping $S$.
\begin{claim}\label{cl:u1}
	Let $x$ and $x+\De$ two points in $\RR^n$. Suppose that
	\begin{enumerate}
		\item
		$S(x)=S(x+\De)=:z$; and
		\item
		for all $i$, $\{x_i+\la \De_i\} \neq z$, $\forall \la\in [0,1]$.
	\end{enumerate}
Then the points $x,x+\De$ fall in the same cell in the tiling induced by $D$.
\end{claim}
\begin{proof}
Suppose towards contradiction that the conclusion of the statement does not hold, i.e.\
$x$ and $x+\De$ belong to different cells in the tiling induced by $D$.
Thus, the rounding function $R$ when applied on $x$ and on $x+\De$ should produce
different lattice points, so there is an $i$ such that $R(x)_i \neq R(x+\De)_i$. We fix that $i$ and
assume without loss of generality that $\De_i\ge 0$ and that $x_i\in [0,1)$. We now consider two cases,
depending on the range $x_i$ falls into:
\begin{enumerate}
  \item If $x_i \in [0,z)$, then by definition of $R$ we get that
  $R(x)_i=0$, and $R(x+\De)_i=0$ unless $x_i+\De_i >z$, which leads to a contradiction to the second condition ($z$ is on the interval between $x_i$ and $x_i+\De_i$).
  \item 	If $x_i \in (z,1)$, then $R(x)_i=1$, and  $R(x+\De)_i=1$ unless $x_i+\De_i >1+z$, which again leads  to a contradiction to the second condition ($1+z$ is on the interval between $x_i$ and $x_i+\De_i$).\qedhere
\end{enumerate}
\end{proof}

Our main technical statement is the following lemma.
\begin{lemma}\label{lem:main_ub}
  There exists a distribution over mappings $(S_{\vec{r}})_{\vec{r}}$ ($\vec{r}$ is a vector
  of randomness) such that for small enough $\eps>0$, setting $\sigma = \eps\frac{\sqrt{\log n}}{n}$ we have
\[
\Expect{\vec{r}}{\Prob{{\bf x}, \bm{\De}\sim\mathcal{N}(0,\sigma^2 I_n)}{\text{Conditions of Claim~\ref{cl:u1} hold for ${\bf x}$ and ${\bf x}+\bm{\De}$}}}
\geq 1 - O(\eps).
\]
\end{lemma}
Deducing from Theorem~\ref{thm:ub} from Lemma~\ref{lem:main_ub} mostly involves measure-theoretic arguments,
and we defer this deduction to Section~\ref{sec:ns_sa_reduce}. We will actually need the following slightly
more informative version of Lemma~\ref{lem:main_ub} above, using the reduction from
mappings to tilings presented in the beginning of this section, and an inspection of the
bodies $D_{\vec{r}}$ our proof gives.

\begin{lemma}\label{lem:main_ub_use}
  There exists a distribution over tiling bodies $(D_{\vec{r}})_{\vec{r}}$ such that

  \begin{enumerate}
    \item For small enough $\eps>0$, we have
    \[
    \Expect{\vec{r}}{\Prob{{\bf x}, \bm{\De}\sim\mathcal{N}(0,\eps^2 I_n)}{\text{At least one of the conditions of Claim~\ref{cl:u1} fail for ${\bf x}$ and ${\bf x}+\bm{\De}$}}}
    \ll \frac{n}{\sqrt{\log n}} \eps.
    \]

    \item For each $\vec{r}$, $D_{\vec{r}}$ is a countable union of semi-algebraic sets (i.e., sets defined by finitely many polynomial inequalities).
  \end{enumerate}
\end{lemma}


\subsection{The construction of $S_{\vec{r}}$}
\subsubsection{Overview}
Before jumping into the technical details, we start with some intuition. Recall that on input $x$ (a set of $n$ points from $\RR/\ZZ$)
we must output a number $z\in\RR/\ZZ$, and our goal is to minimize the probability so that the conditions of Claim~\ref{cl:u1} fail
on a short needle $\ell_{x,\De}$. Note that it would not be beneficial for us to choose $z$ that are close to $x_i$. For example,
if we chose $z$ such that $\card{x_i-z_i}\leq \sigma$, then there is constant probability that the interval $\set{x_i + \lambda \De_i}_{\lambda\in[0,1]}$
would contain the point $z$, i.e.\ the second condition of Claim~\ref{cl:u1} would fail.

Thus, a natural candidate for the choice of $z$ would be the one that maximizes $\min_{i\in[n]}{\card{x_i - z_i}}$. It is not hard to see that this minimum
is typically of the order $\log n/n$, so intuitively the second condition of Claim~\ref{cl:u1} should hold with probability $\geq 1-\eps$. However, such choice
for $z$ would not be very stable: it is typically the case that there are numerous $z_1,\ldots,z_r$ that nearly achieve this maximum, thus the maximizer among
them could change when looking at $x+\De$ (i.e., this event would happen with probability significantly more than $\eps$), leading to a failure of the first
condition of Claim~\ref{cl:u1}.

We must therefore assign each one of the near-maximizer $z_1,\ldots,z_r$ some weight, so that the weight of each one of them does not significantly change when
moving to $x+\De$. A general form of construction of this type is to design a scoring function $f\colon[0,\infty]\to[0,1]$, and given an input $x$ to assign
the weight $w(z) = \prod\limits_{i}{f(\card{x_i - z})}$ to each $z$, and sample $z$ with probability proportional to $w(z)$.

We remark that this general recipe essentially captures our (natural) attempts so far. On the one hand, we want $f$ to penalize $z$
if it is very close to $x_i$, hence we want $f(t)$ at least mildly increasing. On the other hand, if $f$ is very sharply increasing (e.g exponential),
then one runs into the same problems as we had when we thought of picking $z$ that maximizes $\min_{i\in[n]}{\card{x_i - z_i}}$. We are thus
led to consider ``mildly increasing'' scoring functions $f$, and polynomials turn out to be good choice. Indeed, our scoring function $f$ will
be ``trivial'' if $\card{x_i - z}$ is too small or too large (i.e.\ it'll be $0$ if $\card{x_i-z}\leq \frac{\log n}{50 n}$ and $1$ if
$\card{x_i-z}\geq \frac{\log n}{25 n}$), and otherwise behaves cubically.

\subsubsection{A basic scoring function}
Our construction of $(S_{\vec{r}})_{\vec{r}}$ uses a non-negative scoring function $f$ with the following properties.
\begin{fact}\label{fact:std_fn}
There exists a function $f\colon [0,\infty)\to [0,1]$ that is twice differentiable with
continuous second derivative with the following properties:
\begin{enumerate}
  \item $f(t) = 0$ if $t\leq 1$.
  \item $f(t) = 1$ if $t\geq 2$.
  \item $f(t) \ee (t-1)^3$ if $1\leq t\leq 2$.
  \item $\card{f'(t)}\ll t^2$ and $\card{f''(t)}\ll t$ for all $t$.
\end{enumerate}
\end{fact}
Exhibiting function $f$ as in Fact~\ref{fact:std_fn} is not hard, and we omit the proof.
The function $f$ defined by $f(t) = (t-1)^3$ if $1\leq t\leq 2$ and $f(t) = 0$ for $t\leq 1$,
$f(t) = 1$ for $t\geq 2$ is almost enough, except that it is not differentiable at $t=1$. One
can fix by convolving a smooth bump function with compact support.


Next, we wish to define the mapping $S_{\vec{r}}$. We view the input $x$ as a multi-set, and
the randomness vector $\vec{r}$ as an infinite sequence of $(i,h)$ where $i$ is a uniformly
random element from $[m]$ and $h$ is a uniform real-number from $[0,1]$.

Set $m = n^{1/3}$, partition the circle
the circle $\RR/\ZZ$ into $m$ intervals of length $1/m$ each, $I_j:= \left[\frac{j-1}{m},\frac{j}{m}\right]$, and
let $z_j = \frac{j-1/2}{m}$ be the middle of $I_j$.
It will be convenient for us to define $g_j(t) = f(\frac{50n}{\log n}\card{t-z_j})$,
and subsequently $r_j(x):= \prod\limits_{y\in I_j\cap x} g_j(y)$.
There two cases:

\paragraph{Case (A): $r_i(x)\neq 0$ for some $i\in [m]$.}
In this case, we define a probability distribution $p_i(x)$ over the $i$'s proportionally to the $r_i(x)$'s, i.e.\
we define $p_i(x):=\frac{r_i(x)}{\sum_i r_i(x)}$. We now perform correlated sampling of $i\in[m]$ according to $p_i(x)$ using the randomness vector $\vec{r}$.
More precisely, we go over the randomness vector $\vec{r} = (i_1,h_1),(i_2,h_2),\ldots$ and find the smallest $j$ such that
$h_j\leq p_{i_j}(x)$, in which case we choose $i = i_j$. We then define $S_{\vec{r}}(x) = z_{i_j}$.


\paragraph{Case (B): $r_i(x) = 0$ for all $i\in [m]$.}
If $1/2\not\in x$, we define $S_{\vec{r}}(x) = 1/2$.
Otherwise, we define $S_r(x)=z$, where $z$ is the first element from $\{\frac{1}{4n}, \frac{3}{4n},\ldots,\frac{4n-1}{4n}\}$ that is at least $\frac{1}{4n}$-away from all the
entries of $x$.

\subsection{Estimating $g_j$ on close points}
\begin{fact}\label{fact:first_order_g}
Let $j\in[m]$ and
$x_i\in [z_j - \frac{\log n}{25 n}-\eps^{0.95},z_j+\frac{\log n}{25 n}+\eps^{0.95}]\setminus [z_j - \frac{\log n}{50 n},z_j+\frac{\log n}{50 n}]$,
$\De_i\in\mathbb{R}$, and denote
$\alpha_i = {\sf dist}\left(x_i, [z_j - \frac{\log n}{50 n}, z_j+\frac{\log n}{50n}]\right)$.

\begin{enumerate}
  \item If $\alpha_i\geq 2\card{\De_i}$, then $\card{g_j(x_i+\De_i) - g_j(x_i)}\ll \frac{\card{\De_i}}{\alpha_i} g_j(x_i)$.
  \item In general, $\card{g_j(x_i+\De_i) - g_j(x_i)}\ll n^3(\alpha_i^3+\card{\De_i}^3)$.
\end{enumerate}

\end{fact}
\begin{proof}
     Using
     Taylor's approximation with remainder, there is $y_i\in [x_i,x_i + \De_i]$ such that
     $g_j(x_i+\De_i) = g_j(x_i) + g_j'(y_i)\De_i$, hence
     \begin{align*}
     \card{g_j(x_i+\De_i) - g_j(x_i)}
     \ll \card{\De_i} \card{g_j'(y_i)}
     &\ll \card{\De_i}\frac{50n}{\log n}f'\left(\frac{50n}{\log n}\card{y_i-z_j}\right)\\
     &\ll \frac{n}{\log n} \card{\De_i} \left(\frac{50n}{\log n}\card{y_i-z_j}-1\right)^2.
     \end{align*}
     For the second item, since $y_i\in [x_i,x_i+\De_i]$,
     we get that $\card{\frac{50n}{\log n}\card{y_i-z_j}-1}\leq \frac{50n}{\log n}\left(\alpha_i + \card{\De_i}\right)$,
     and plugging that in yields
     \[
     \card{g_j(x_i+\De_i) - g_j(x_i)}
     \ll n^3\card{\De_i}(\alpha_i^2 + \De_i^2)
     \ll n^3(\alpha_i^3+\card{\De_i}^3),
     \]
     where the last inequality holds as $ab\ll a^3 + b^{3/2}$ for all $a,b>0$ (Young's inequality).
     For the first item, note that since $y_i\in [x_i,x_i+\De_i]$ we get that
     $\left(\frac{50n}{\log n}\card{y_i-z_j}-1\right)\geq \frac{50n}{\log n}\left(\alpha_i - \card{\De_i}\right)$,
     and by the lower bound on $\alpha_i$ this is $\geq \frac{25n}{\log n}\alpha_i$. Therefore we may continue
     as
     \[
     \card{g_j(x_i+\De_i) - g_j(x_i)}
     \ll \frac{n}{\log n} \card{\De_i} \left(\frac{50n}{\log n}\card{y_i-z_j}-1\right)^2
     \ll \frac{\card{\De_i}}{\alpha_i}\left(\frac{50n}{\log n}\card{y_i-z_j}-1\right)^3.
     \]
     Also, we have that
     $\left(\frac{50n}{\log n}\card{y_i-z_j}-1\right)\leq \frac{50n}{\log n}\left(\alpha_i + \card{\De_i}\right)
     \ll \frac{n}{\log n} \alpha_i$, so
     \[
     \card{g_j(x_i+\De_i) - g_j(x_i)}
     \ll \frac{\card{\De_i}}{\alpha_i}\left(\frac{n}{\log n} \alpha_i\right)^3
     \ll \frac{\card{\De_i}}{\alpha_i}g_j(x_i).\qedhere
     \]

\end{proof}
\subsection{Analysis of the construction}\label{sec:analysis_rare}
In this section we prove that Lemma~\ref{lem:main_ub} holds for the construction of $S_{\vec{r}}$ from the last section, and for
that we show that for small enough $\eps$, the expected probability of the complement event is $O(\eps)$, i.e.\ that
\begin{equation}\label{eq:main_ub}
\Expect{\vec{r}}{\Prob{{\bf x}, \bm{\De}\sim\mathcal{N}(0, \sigma^2 I_n)}
{\text{One of the conditions in Claim~\ref{cl:u1} fails for ${\bf x}$ and ${\bf x}+\bm{\De}$}}}
\ll \eps.
\end{equation}
We will think of $\eps$ as very small (say $\eps \leq 2^{-n^2}$), and analyze the contribution of $x$'s from case (A) and case (B) separately.
Case (A) is the main case that occurs often, and case (B) should be thought of rare.


\subsubsection{Analysis of case (B)}
First, we show that the probability ${\bf x}$ (or equivalently ${\bf x}+\bm{\De}$) falls into Case~(B) is at most $n^{-\omega(1)}$.
For this, it will be helpful for us to sample ${\bf x}$, a multi-set of $n$ uniformly chosen numbers in $[0,1]$ in the following equivalent way:
\begin{itemize}
	\item
	Sample ${\bf t}_1,\ldots,{\bf t}_m$ --- where ${\bf t}_i$ is the number of $i$'s such that ${\bf x}_i$'s that fall into interval $I_i$.
	\item Sample ${\bf t}_i$ points uniformly from $I_i$, for each $i=1,\ldots,m$.
\end{itemize}
Note that $\E[{\bf t}_i]=n/m$, hence by Chernoff bound $\Pr[{\bf t}_i\ge 2 \cdot n/m] =e^{-\Omega(n/m)}=n^{-\omega(1)}$. Thus, by
the union bound we have that
\[
\Prob{}{\forall i ~{\bf t}_i\leq 2\cdot n/m} = 1 - n\cdot n^{-\omega(1)} = 1-n^{-\omega(1)}.
\]
Next, we condition on ${\bf t}_i = t_i$, and assume that indeed $t_i<2\cdot n/m$ for all $i$. Let $E_i$ be the event that $r_i(x)=0$.
Note that conditioned on ${\bf t}_i = t_i$, the $E_i$'s are independent and that
\begin{align}
\Pr[\neg E_i|~t_1,\ldots,t_m]
=\Prob{{\bf a}\in I_i}{\frac{50 n}{\log n}\card{{\bf a}-z_i}\leq 1}^{t_i}
=\left( 1-\frac{\log n/25n}{1/m} \right)^{t_i}
&\geq
\left( 1-\frac{m \log n}{25 n} \right)^{2 \cdot n/m}\notag\\\label{eq:notEi}
&\geq e^{-4\log n/25}
= n^{-4/25},
\end{align}
where we used the fact that $e^{-\delta}\leq 1-\delta/2$ for small enough $\delta>0$.
Therefore,
$$
\Pr[E_1\wedge E_2 \wedge \ldots \wedge E_m|~t_1,\ldots,t_m] 
\leq (1-n^{-4/25})^m
=(1-n^{-4/25})^{n^{1/3}} = n^{-\omega(1)},
$$
as long as the $t_i$'s satisfy the condition $t_i<2\cdot n/m$. Therefore, the overall probability of case~(B) is $n^{-\omega(1)}$.

Next, we analyze the probability that the conditions of Claim~\ref{cl:u1} fail given we are in case (B).
Note that if the conditions of Claim~\ref{cl:u1} fail to hold, then either (I) exactly one of ${\bf x}$, ${\bf x}+\bm{\De}$ falls under Case~(B),
or (II) both ${\bf x}$ and ${\bf x}+\bm{\De}$ fall under Case~(B),  but $1/2\in {\bf x}+\la \bm{\De}$ for some $\la\in [0,1]$. We'll bound these cases separately.



\paragraph{Case (I).} Assuming ${\bf x}$ is under Case~(B), we know that each
of the $m$ intervals of the form $J_i := [z_i-\frac{\log n}{50n},z_i+\frac{\log n}{50n}]$ contains at least one point from ${\bf x}$.
Let ${\bf x}_i$ be that point (if there are multiple, pick one at random). Then ${\bf x}_i$ is uniformly distributed in $J_i$.
Therefore, the probability of ${\bf x}_i+\bm{\De}_i$ is outside $J_i$, where
$\bm{\De}_i \sim N(0,\sigma^2)$ and $\sigma^2\leq \eps^2$, is $O(\eps)$.
Given case (B) occurs with probability $\leq n^{-\omega(1)}$, we conclude that its contribution
to the conditions of Claim~\ref{cl:u1} failing is at most
\[
O(m \eps)\cdot n^{-\omega(1)} = O(\eps).
\]

\paragraph{Case (II).}
Fix $\bm{\De} = \De$, and consider ${\bf x}_j$ conditioned on being in case (B).
If ${\bf x}_j$ is in one of the intervals $J_i$, then its distribution is uniform over $J_i$, in which
case we get that the probability $1/2$ falls inside the interval $[{\bf x}_j,{\bf x}_j + \De_j]$ is at most
$m\card{\De_j}$. If ${\bf x}_j$ is not in one of the intervals $J_i$, then it is distributed uniformly
on $[0,1]\setminus\cup_{i=1}^{m} J_i$, and the probability $1/2$ is in $[{\bf x}_j,{\bf x}_j + \De_j]$ is
at most $2\card{\De_j}\leq m\card{\De_j}$.

Therefore by the union bound,
\[
\cProb{{\bf x}}{{\sf case (B)}, \De}{\exists j\in [n]~1/2\in [{\bf x}_j,{\bf x}_j+\De_j]}
\leq m\sum\limits_{j=1}^{n}{\card{\De_i}}.
\]
Taking expectation over $\De\sim\mathcal{N}(0,\sigma^2 I_n)$ and
using Cauchy-Schwarz we get that
\[
\cProb{{\bf x},\bm{\De}}{{\sf case (B)}}{\exists j\in [n]~~1/2\in [{\bf x}_j,{\bf x}_j+\bm{\De}_j]}
\leq m\Expect{\bm{\De}}{\sum\limits_{j=1}^{n}{\card{\bm{\De}_i}}}
\leq m\sqrt{n}\sqrt{\Expect{\bm{\De}\sim\mathcal{N}(0,\sigma^2 I_n)}{\norm{\bm{\De}}_2^2}}
= mn\sigma.
\]
Therefore, the contribution of this case is upper bounded as
\[
\Prob{{\bf x},\bm{\De}}{{\sf case (B)}\land \exists j\in [n]~1/2\in [{\bf x}_j,{\bf x}_j+\De_j]}
\leq \Prob{{\bf x},\bm{\De}}{{\sf case (B)}}mn\sigma
=n^{-\omega(1)}\cdot \sigma = O(\eps).
\]

\subsubsection{Analysis of case (A)}
We now analyze the contribution of $x$'s that fall into case (A) to the left hand side of~\eqref{eq:main_ub}.
\paragraph{Case~(A), Condition~2.}  If  ${\bf x}$  falls under Case~(A),
then the distance from all ${\bf x}_i$'s to $z=S({\bf x})$ is at least $\frac{\log n}{100 n}$.
Therefore, Condition~2 holds as long as $|\bm{\De}_i| < \frac{\log n}{100 n}$ for all $i$.
Since for each $i$ we have that
\[
\Prob{\bm{\De}\sim\mathcal{N}(0,\sigma^2 I_n)}{\card{\bm{\De}_i}\geq \frac{\log n}{100 n}}
=
\Prob{\bm{\De}\sim\mathcal{N}(0,\sigma^2 I_n)}{\bm{\De}_i^2\geq \frac{\log^2 n}{100^2 n^2}}
\ll \frac{\sigma^2}{\log ^2 n/n^2}
\ll \eps^2/\log n,
\]
we get by the union bound that
\[
\Prob{\bm{\De}\sim\mathcal{N}(0,\sigma^2 I_n)}{\exists i~\card{\bm{\De}_i}\geq \frac{\log n}{100 n}}
\ll n\eps^2/\log n
\ll \eps,
\]
for a sufficiently small $\eps$.


\subsection{Case~(A), Condition~1.}
This is the main part of the proof.
We show that in case (A), the probability that $S_{\vec{r}}({\bf x})\neq S_{\vec{r}}({\bf x}+\bm{\De})$
is at most $O(\eps)$. Note that the procedure describing $S_{\vec{r}}$ in this case
is the correlated sampling  procedure of Holenstein~\cite{Holenstein}, where
$S_{\vec{r}}(x)$ samples $i$ according to
the distribution $p(x) = (p_1(x),\ldots, p_m(x))$ and
$S_{\vec{r}}(x+\De)$ samples $i$ according to the distribution $p(x+\De)$. Therefore,
the probability they sample different $i$'s is at most the statistical distance
between the distributions, $\norm{p(x) - p(x+\De)}_1$. Therefore, we must show that
\begin{equation}\label{eq:mainA1}
\cbExpect{{\bf x},{\bf x}+\bm{\De}}{{\sf case (A)}}{\norm{p({\bf x}) - p({\bf x}+\bm{\De})}_1} = O(\ve).
\end{equation}

Before we turn to this task, we upper bound the contribution from several rare cases.

\subsubsection{Contribution from some rare cases}
First, we show that the case some $\bm{\De}_i$ is too large contributes at most $O(\eps)$
to the LHS of~\eqref{eq:mainA1}.
\begin{claim}\label{cl:A-1}
	$\Prob{\bm{\De}\sim\mathcal{N}(0,\sigma^2 I_n)}{\card{\bm{\De}_i}\geq \eps^{0.95}/n~\text{for some $i$}}\leq \eps$.
\end{claim}
\begin{proof}
For each $i$, we have that
\[
\Prob{\bm{\De}\sim\mathcal{N}(0,\sigma^2 I_n)}{\card{\bm{\De}_i}\geq \eps^{0.95}/n}
\leq 2^{-\Omega((\eps^{0.95}/n)^2/\sigma^2)} = 2^{-\Omega\left(\frac{1}{\eps^{0.1}\log n}\right)}
\leq \frac{\eps}{n},
\]
for small enough $\eps$,
and the claim follows from the union bound.
\end{proof}
From now on, we assume that the $\bm{\De}_i$'s are distributed from $N(0,\sigma^2)|_{|\bm{\De}_i| < \eps^{0.95}/n}$.
In particular, we can assume that if ${\bf t}_j$ is the number of $\bm{x}$'s that fall into interval $I_j$, these numbers stay the same under ${\bf x}+\bm{\De}$.
\footnote{Strictly speaking, $x_i + \De_i$ may be in a different interval than $x_i$, but in this case it doesn't affects the distribution $p(x)$.
Indeed, suppose $x_i$ is in $I_j$ but $x_i + \De_i$ is in $I_{j+1}$. Then
$\card{x_i + \De_i - z_{j+1}}\geq \card{z_{j+1} - j/m} - \card{\De_i} - \card{x_i - j/m}
\geq 1/m - 2\eps^{0.95}/n\geq 1/m - \eps^{0.95}$. Therefore, $\frac{50n}{\log n}\card{x_i + \De_i - z_{j+1}} > 2$, and
so $f(\frac{50n}{\log n}\card{x_i + \De_i - z_{j+1}}) = 1$.
}
Next, we handle the case in which $p({\bf x})$ is supported only on a single $j$.
Note that in this case, if $p({\bf x}+\bm{\De})$ is also only supported on this single $j$,
then the contribution of these cases to the LHS of~\eqref{eq:mainA1} is $0$.
We show that the contribution from the other case is $O(\eps)$.
\begin{claim}\label{cl:A-3}
	\[\Prob{{\bf x},\bm{\De}}{\text{$\exists j^{\star}$ such that $p({\bf x})$ is only supported on $j^{\star}$
        and the support of $p({\bf x}+\bm{\De})$ is different}}\ll\eps.\]
\end{claim}
\begin{proof}
    In case (B), we have shown that the probability that $r_j({\bf x}) = 0$ for all $j$ is $n^{-\omega(1)}$, and the same argument
    shows that the probability $r_j({\bf x}) = 0$ for all but a single $j^{\star}$ is still $n^{-\omega(1)}$. Denote this event by $E$.

    Let us condition on the event $E$, on $j^{\star}$ and the number ${\bf t}_1,\ldots,{\bf t}_m$ of ${\bf x}_i$'s that fall into $I_1,\ldots,I_m$.
    Note that for each $j\neq j^{\star}$, since $r_j({\bf x}) = 0$ there is $i$ such that ${\bf x}_i\in J_j \defeq [z_{j} - \frac{\log n}{50 n}, z_{j} + \frac{\log n}{50 n}]$,
    and we condition on that $i_j$ for each $j$ (if there is more than one, we choose one arbitrarily). Note that the distribution of ${\bf x}_{i_j}$
    is thus uniform over $J_j$.

    Now note that if for each $j\neq j^{\star}$ it holds that ${\bf x}_{i_j} + \bm{\De}_{i_j}\in J_j$, then $r_j({\bf x}+{\bf \De}) = 0$, so the only contribution to
    the probability of the event in question comes when ${\bf x}_{i_j} + \bm{\De}_{i_j}\not\in J_j$ (or from case (B), which we have already accounted for earlier).
    Conditioned on $\bm{\De} = \De$, the probability for that is
    at most
    \[
    \Expect{({\bf x}_{i_j})_{j\neq j^{\star}}}{\sum\limits_{j\neq j^{\star}}1_{{\bf x}_{i_j} + \De_{i_j}\not\in J_j}}
    = \sum\limits_{j\neq j^{\star}}\Expect{{\bf x}_j}{1_{{\bf x}_{i_j} + \De_{i_j}\not\in J_j}}
    \leq \sum\limits_{j\neq j^{\star}}\frac{\card{\De_j}}{\log n/(50n)},
    \]
    therefore taking expectation over $\De$ and using Cauchy-Schwarz we get that
    \[
    \Expect{\bm{\De},({\bf x}_{i_j})_{j\neq j^{\star}}}{\sum\limits_{j\neq j^{\star}}1_{{\bf x}_{i_j} + \bm{\De}_{i_j}\not\in J_j}}
    \ll
    \frac{n}{\log n}\sqrt{m}\sqrt{\sum\limits_{j\neq j^{\star}}
    \Expect{\bm{\De}}{\card{\bm{\De}_j}^2}}
    \leq \frac{n}{\log n}\sqrt{m}\sqrt{m\sigma^2} \leq n^2\sigma.
    \]
    Therefore, we get that
    \[
    \Prob{{\bf x},\bm{\De}}{\text{$p({\bf x})$ is only supported on $j^{\star}$,
        but the support of $p({\bf x}+\bm{\De})$ is different}}
        \leq \Prob{}{E}n^2\sigma
        \leq n^{-\omega(1)} n^2 \sigma
        \ll \eps.\qedhere
    \]
\end{proof}
Let $E$ be the event that the support of $p({\bf x})$ consists of at least two distinct $j$'s. We condition on the event $E$ in the subsequent argument.
The following claim shows that conditioned on $E$, the sum of the $r_j({\bf x})$'s is at least somewhat bounded away from $0$. It will only come into play
later in the proof.
\begin{claim}\label{cl:A-2}
  $\cProb{{\bf x}}{E}{\sum\limits_{j} r_j({\bf x}) \leq \eps^{1.6}} \ll \eps$.
\end{claim}
\begin{proof}
Since we conditioned on $E$, there are $j_1\neq j_2$'s such that
$r_{j_1}({\bf x}),r_{j_2}({\bf x}) > 0$. We condition on $j_1$ and $j_2$, and assume without loss of generality that
$j_1 = 1,j_2=2$. We show that
\begin{equation}\label{eq:cl:A-2:1}
\cProb{{\bf x}}{r_1({\bf x}),r_2({\bf x})>0}{r_{1}({\bf x})<\ve^{1.6} \wedge r_{2}({\bf x})<\ve^{1.6}} \ll \eps,
\end{equation}
and thus the result would follow.

Let ${\bf t}_1$ be the number of $i$'s such that ${\bf x}_i\in I_{1}$, and ${\bf t}_2$ be the number of $i$'s such that ${\bf x}_i\in I_{2}$.
Note that ${\bf t}_1, {\bf t}_2 \leq n$.
In addition, conditioned on ${\bf t}_1=t_1$ and ${\bf t}_2=t_2$, the events $r_1({\bf x})<\ve^{1.6}$ and $r_2(x)<\ve^{1.6}$ become independent.
Therefore, to prove \eqref{eq:cl:A-2:1}, it suffices to show for all $t_1\leq n, t_2\leq n$,
  	\begin{equation}
  \cProb{{\bf x}}{r_1({\bf x})>0,~{\bf t}_1=t_1,~{\bf t}_2=t_2}{r_1({\bf x})<\ve^{1.6}}\ll \eps^{0.5}.
  \label{eq:cl:A-2:2}
  \end{equation}
  Note that one way to sample $r_1({\bf x})| r_1({\bf x})>0,~{\bf t}_1=t_1$ is as follows.
  \begin{itemize}
  	\item Sample points ${\bf x}_1,\ldots,{\bf x}_{t_1}$ uniformly from $I_1$ conditioned on $|{\bf x}_i-z_1| > \frac{\log n}{50 n}$;
  	\item
  	$r_1({\bf x}) = \prod_{i=1}^{t_1} g_1({\bf x}_i)$.
  	\end{itemize}
  Let ${\bf Y}_i$ be the random variable ${\bf Y}_i := g_1({\bf x}_i)^{-0.32}$, where
  ${\bf x}_i$ is sampled as above (we need $0.32<1/3$).
  Let $E$ be the event that $\card{{\bf x}_i - z_1}\geq \frac{\log n}{25 n}$.
  If $E$ holds, then we get that $g_1({\bf x}_i) = 1$, and
  otherwise $g_1({\bf x}_i)\gg \card{\frac{50n}{\log n}\card{{\bf x}_i - z_1} - 1}^3$,
  so
  \[
  \Expect{}{{\bf Y}_i}
  \leq \Prob{}{E}\cdot 1 + \Prob{}{\bar{E}}\cExpect{}{\bar{E}}{g_1({\bf x}_i)^{-0.32}}
  \ll 1 + \cExpect{}{\bar{E}}{\card{\frac{50n}{\log n}\card{{\bf x}_i - z_1} - 1}^{-0.96}}.
  \]
  We write the last expectation as an integral, noting that $\card{{\bf x}_i-z_1}$ is distributed uniformly on
  $\left[\frac{\log n}{50 n},\frac{\log n}{25 n}\right]$, hence
  \[
  \cExpect{}{\bar{E}}{\card{\frac{50n}{\log n}\card{{\bf x}_i - z_1} - 1}^{-0.96}}
  \ll\frac{n}{\log n}\int_{\frac{\log n}{50n}}^{\frac{\log n}{25n}}{\card{\frac{50n}{\log n}t - 1}^{-0.96} dt}
  =\frac{1}{50}\int_{0}^{1}{y^{-0.96} dt}
  \ll 1,
  \]
  where we made the change of variables $y = \frac{50n}{\log n}t - 1$.
  Thus, $\E[{\bf Y}_i] \ll 1$, and so there is a constant $B$ such that $\E[{\bf Y}_i]\leq B$.
  Therefore by independence $\E\left[\prod_{i=1}^{t_1} {\bf Y}_i\right]\leq B^{t_1}\leq B^n$,
  and so writing $r_1({\bf x})$ in terms of the ${\bf Y}_i$'s and using Markov's inequality we get that
  \[
    \cProb{{\bf x}}{r_1({\bf x})>0,~{\bf t}_1=t_1,~{\bf t}_2=t_2}{r_1({\bf x})<\ve^{1.6}}
    = \Pr\left[\prod_{i=1}^{t_1} {\bf Y}_i > \ve^{-1.6\times 0.32} \right] \leq B^n\cdot \ve^{0.512} \ll \eps^{0.5}.
  \qedhere
  \]
\end{proof}

\subsubsection{Analyzing the typical case}
To expand out $\norm{p({\bf x}) - p({\bf x}+\bm{\De})}_1$, we will be using the following claim. The set-up one should
have in mind is that $r_j = r_j(x)$ and $d_j = r_j(x+\De)$ for some $x$ and $\De$ that are typical enough.
\begin{claim} \label{cl:A-6}
	Let $r_j\geq 0$, $d_j$ be real-numbers satisfying $\card{d_j}\leq r_j/2$ for all $j$.
    Denote $T=\sum r_j$, $T'=\sum (r_j+d_j)$, and let $p_j = r_j/T$ and $q_j = (r_i+d_i)/T'$ be two distributions.
    Then
	\begin{equation}\label{eq:cl:A-6:0}
	\norm{p-q}_1 \ll \sum_i \frac{|d_i|}{r_i} \cdot \frac{\min(r_i, T-r_i)}{T}.
	\end{equation}
\end{claim}
We defer the proof of Claim~\ref{cl:A-6} to Section~\ref{sec:def_proofs}. Morally speaking, it says that
\begin{equation}\label{eq:pretend1}
\Expect{{\bf x},\bm{\De}}{\norm{p({\bf x}) - p({\bf x}+\bm{\De})}_1}
\ll \sum\limits_{j=1}^{m}{\Expect{{\bf x}}{\Expect{\bm{\De}}{\frac{\card{r_j({\bf x}) - r_j({\bf x}+\bm{\De})}}{r_j({\bf x})}\cdot\frac{\min(r_j({\bf x}),T({\bf x})-r_j({\bf x}))}{T({\bf x})}}}},
\end{equation}
where $T(x) = \sum\limits_{j}{r_j(x)}$ (this is only morally because we are assuming that the supports of $p_j(x)$ and $p_j(x+\De)$
are the same, but formally speaking they may be different). In particular, to be able to handle with that we first must understand
the expectation of $\card{r_j(x) - r_j(x+\bm{\De})}$ over $\bm{\De}$.
\begin{claim}\label{cl:A-4}
    Let $j\in [m]$,  $x_1,\ldots,x_k \in [z_j-\frac{\log n}{25 n}-\eps^{0.95},z_j+\frac{\log n}{25 n}+\eps^{0.95}] \setminus [z_j-\frac{\log n}{50 n},z_j+\frac{\log n}{50 n}]$,
    and let $r(x)=\prod_{i=1}^c g_j(x_i)$.
	Denote $\al_i =\text{dist}\left(x_i, [z_j-\frac{\log n}{50n},z_j+\frac{\log n}{50n}]\right)$.
	and let $\bm{\De}_i \sim N(0,\sigma^2)|_{|\De_i|<\ve^{0.95}}$.
	Then
	\begin{equation}
	\E_{\bm{\De}}[|r(x+\bm{\De})-r(x)|] \ll \max\left(\eps^{2.65},r(x)\cdot \sigma\cdot\sqrt{\sum_{i=1}^c \frac{1}{\al_i^2}}\right).
	\end{equation}
\end{claim}

\begin{proof}
	We consider two cases.
	
	\paragraph{Case 1: $\al_i\le \ve^{0.9}$ for some $i$.}
    In this case, we have
    \[
    g_j(x_i)\ll
    \left(\frac{50n}{\log n}\alpha_i\right)^3
    \ll n^3\eps^{3\cdot 0.9}
    \ll \eps^{2.66}.
    \]
    Similarly, we have ${\sf dist}(x_i+\bm{\De}_i,[z_j-\frac{\log n}{50n},z_j+\frac{\log n}{50}])
    \leq \alpha_i + \card{\bm{\De}_i}\leq 2\alpha_i$, so $g_j(x_i+\bm{\De}_i)\ll \eps^{2.66}$.
    We conclude that $r(x_i),r(x_i+\bm{\De}_i)\ll \eps^{2.66}$, hence the contribution from these cases is at most $\eps^{2.65}$.
	
	 \paragraph{Case 2: $\al_i>\ve^{0.9}$ for all $i$.} In this case, we get that $x_i+\bm{\De}_i$ is also not in
     the interval $[z_j-\frac{\log n}{50n},z_j+\frac{\log n}{50n}]$, hence $g_j(x_i+\bm{\De}_i)\neq 0$,
     so $r(x+\bm{\De})>0$. Since $r(x)$ are defined using products, it would be more convenient for us to analyze
     $\log(r(x+\bm{\De})/r(x))$ as opposed to $r(x+\bm{\De})/r(x)-1$, and to justify we can do that we first argue that
     $r(x+\bm{\De})/r(x) = 1+o(1)$.

     To see that, note that as $\card{\bm{\De}_i}\leq \eps^{0.95}\leq \alpha_i/2$, we may use Fact ~\ref{fact:first_order_g} to conclude that
     \[
     \card{g(x_i+\bm{\De}_i) - g(x_i)}
     \ll \frac{\card{\bm{\De}_i}}{\alpha_i} \card{g(x_i)}
     \ll \eps^{0.05} \card{g(x_i)}.
     \]
     In particular, we get that
     $\frac{g_j(x_i+\bm{\De}_i)}{g_j(x_i)} = 1\pm O(\eps^{0.05})$, and hence
     $\frac{r(x+\bm{\De})}{r(x)} = 1\pm O(k\eps^{0.05})$. Writing $\frac{r(x+\bm{\De})}{r(x)} = 1+\bm{\eta}$,
     we get $\bm{\eta}$ is small in absolute value, and hence
     $\card{\log(r(x+\bm{\De})/r(x))}\gg \card{\bm{\eta}}\gg \card{\frac{r(x+\bm{\De})}{r(x)} - 1} = \card{\frac{r(x+\bm{\De})-r(x)}{r(x)}}$.
     I.e.,
	 \begin{equation}
	 \Expect{\bm{\De}}{\frac{|r(x+\bm{\De})-r(x)|}{r(x)}}
     \ll \Expect{\bm{\De}}{\card{\log \left(\frac{r(x+\bm{\De})}{r(x)}\right)}}
     = \Expect{\bm{\De}}{\left|\sum_{i=1}^k \log \frac{ g_j(x_i+\bm{\De}_i)}{g_j(x_i)} \right|}
     = \Expect{\bm{\De}}{\left|\sum_{i=1}^k {\bf Y}_i \right|},
	 	  \label{eq:cl:A-4:1}
\end{equation}
	 where we define the random variables ${\bf Y}_i = \log \frac{ g_j(x_i+\bm{\De}_i)}{g_j(x_i)}$.

     Observe that ${\bf Y}_i$'s are mutually independent, since each ${\bf Y}_i$ only depends on
	 the corresponding $\bm{\De}_i$. We wish to upper bound the average and variance of ${\bf Y}_i$,
     and to do that it would be more convenient to analyze ${\bf Z}_i = \frac{ g_j(x_i+\bm{\De}_i)-g_j(x_i)}{g_j(x_i)}$
     and then relate the two.

     Using second order Taylor's approximation, we have that there is ${\bf y}_i\in[x_i,x_i+{\bf \De}_i]$ such that
     \[
       g_j(x_i+{\bf \De}_i) = g_j(x_i) + g_j'(x_i){\bf \De}_i + \half g_j''({\bf y}_i){\bf \De}_i^2,
     \]
     hence
     \begin{equation}\label{eq:avg_Z}
       \card{\Expect{{\bf \De}}{{\bf Z}_i}}
       =\frac{1}{g_j(x_i)}\card{\Expect{\De}{g_j'(x_i){\bf \De}_i + \half g_j''({\bf y}_i){\bf \De}_i^2}}
       =\frac{1}{2g_j(x_i)}\card{\Expect{\De}{g_j''({\bf y}_i){\bf \De}_i^2}}.
     \end{equation}
     Using properties of $f$, we have
     \[
     \card{g_j''({\bf y}_i)} =
     \left(\frac{50n}{\log n}\right)^2 f''\left(\frac{50n}{\log n}\card{{\bf y}_i-z_j}\right)
     \ll \left(\frac{50n}{\log n}\right)^2 \card{\frac{50n}{\log n}\card{{\bf y}_i-z_j}-1}.
     \]
     Since ${\bf y}_i\in [x_i,x_i+{\bf \De}_i]$, we get that
     $
     \left(\frac{50n}{\log n}\card{{\bf y}_i-z_j}-1\right)
     \geq \frac{50n}{\log n}\alpha_i - \eps^{0.95}
     \geq \frac{25n}{\log n}\alpha_i
     $,
     and so we may continue the previous inequality as
     \[
     \card{g''({\bf y}_i)}
     \ll \left(\frac{50n}{\log n}\right)^2\frac{\card{\frac{50n}{\log n}\card{{\bf y}_i-z_j}-1}^3}{(\frac{25n}{\log n}\alpha_i)^2}
     \ll \frac{1}{\alpha_i^2} \card{g_j({\bf y}_i)}
     \ll \frac{1}{\alpha_i^2} \card{g_j(x_i)},
     \]
     where the last inequality is by Fact~\ref{fact:first_order_g}.
     Plugging this into~\eqref{eq:avg_Z} we get that
     \[
     \card{\Expect{\bm{\De}}{{\bf Z}_i}}
     \ll \frac{1}{\alpha_i^2} \Expect{\bm{\De}}{\bm{\De}_i^2}
     =\frac{1}{\alpha_i^2} \sigma^2.
     \]

	 In a similar fashion, we upper bound the second moment of ${\bf Z}_i$.
     Using Fact~\ref{fact:first_order_g}, we get that $\card{{\bf Z}_i}\leq \frac{\bm{\De}_i}{\alpha_i}$, and
     so $\Expect{\bm{\De}}{{\bf Z}_i^2}\ll \frac{1}{\alpha_i^2}\Expect{\bm{\De}}{\bm{\De}_i^2} =\frac{1}{\alpha_i^2} \sigma^2$.




     We can now upper bound the average of ${\bf Y}_i$ as follows. Recall that, $\card{{\bf Z}_i} = o(1)$ so by Taylor's approximation
     ${\bf Y}_i = \log(1+{\bf Z}_i) = {\bf Z}_i - \frac{1}{2(1+\bm{\xi}_i)^2} {\bf Z}_i^2$ for some $\bm{\xi}_i\in [1,1+{\bf Z}_i]$ and hence
	  \begin{equation}
	 \card{\Expect{}{{\bf Y}_i}} \ll \card{\E[{\bf Z}_i]} + \card{\E[{\bf Z}_i^2]|} \ll \frac{1}{\alpha_i^2}\sigma^2.
	 \label{eq:cl:A-4:4}
	 \end{equation}
	 This approximation (along with the fact that $\card{{\bf Z}_i} = o(1)$) also implies $\card{{\bf Y}_i}\ll \card{{\bf Z}_i}$, hence
	 \begin{equation}
	 \E[{\bf Y}_i^2] \ll \E[{\bf Z}_i^2]  \ll \frac{1}{\alpha_i^2}\sigma^2.
	 \label{eq:cl:A-4:5}
	 \end{equation}
	
	 We can now continue equation~\eqref{eq:cl:A-4:1} to upper bound the LHS there.
     Denoting $\mu_i:=\E[{\bf Y}_i]$, we have
	 \begin{align*}
	 \Expect{\De}{\card{\sum_{i=1}^k {\bf Y}_i}}
     \leq \sum_{i=1}^k |\mu_i|  +  \Expect{\De}{\card{\sum_{i=1}^k {\bf Y}_i-\mu_i}}
     \le \sum_{i=1}^k |\mu_i|
	 + \sqrt{\Expect{\De}{\sum_{i=1}^k ({\bf Y}_i-\mu_i)^2}},
	 \end{align*}
     where in the last inequality we used Cauchy-Schwarz and the fact that $Y_i$'s are independent.
     Using ~\eqref{eq:cl:A-4:4} we have that $\sum_{i=1}^k |\mu_i|\ll \sigma^2\sum_{i=1}^k \frac{1}{\alpha_i^2}$,
     and to upper bound the second term we use ~\eqref{eq:cl:A-4:5}:
     \[
     \Expect{\bm{\De}}{({\bf Y}_i-\mu_i)^2}
     \leq \Expect{\bm{\De}}{{\bf Y}_i^2}
     \ll \frac{1}{\alpha_i^2}\sigma^2.
     \]
     Together, we get that
     \[
     \Expect{\bm{\De}}{\card{\sum_{i=1}^k {\bf Y}_i}}
     \ll \sigma^2\sum_{i=1}^k \frac{1}{\alpha_i^2} + \sqrt{\sigma^2\sum_{i=1}^k \frac{1}{\alpha_i^2}}
     \ll \sigma  \sqrt{\sum_{i=1}^k \frac{1}{\alpha_i^2}},
     \]
     where the last inequality holds since
     $\sigma^2\sum_{i=1}^k \frac{1}{\alpha_i^2}\ll 1$
     (as $\sigma^2\ll \eps^2$ and $\alpha_i\geq \eps^{0.9}$).
\end{proof}


Next, using the previous claim we upper bound the expectation of each summand on the RHS of~\eqref{eq:pretend1}.
The following statement addresses a single term, and should be thought of as being applied after conditioning on
$x,\De$ being not-too untypical, and focusing only on $x_i$'s for which there is a chance that
$g_j(x_i + \De_i)\neq g_j(x_i)$.

\begin{claim}\label{cl:A-5}
Let $j\in [m]$, $k\leq n$, $S\geq 0$ and let $x_1,\ldots,x_k$
be chosen uniformly at random from
$[z_j-\frac{\log n}{25 n}-\eps^{0.95},z_j+\frac{\log n}{25 n}+\eps^{0.95}] \setminus [z_j-\frac{\log n}{50 n},z_j+\frac{\log n}{50 n}]$.
Let $\De_i \sim N(0,\sigma^2)|_{|\De_i|<\ve^{0.95}}$.
Then
\begin{align*}
\E_{{\bf x},\bm{\De}} \left[ \frac{|r_j({\bf x}+\bm{\De})-r_j({\bf x})|}{r_j({\bf x})}\cdot \frac{\min(r_j({\bf x}),S)}
 {r_j({\bf x})+S+\ve^{1.6}}\right]
\ll  \ve^{1.05}+ k\frac{\sigma n}{\log n}  \cdot \Pr_{{\bf x}}[r_j({\bf x})\ge S]\\
+\sigma \frac{n}{\log n}\sqrt{k} \Expect{{\bf x}}{\frac{r_j({\bf x})}{r_j({\bf x})+S}}.
\end{align*}
\end{claim}
\begin{proof}
	Upper bounding $\max(a,b)\leq a+b$ for $a,b\geq 0$, by Claim~\ref{cl:A-4}, we have
	\begin{align*}
	\E_{{\bf x},\bm{\De}} \left[ \frac{|r_j({\bf x}+\bm{\De})-r_j({\bf x})|}{r_j({\bf x})}\cdot \frac{\min(r_j({\bf x}),S)}
	{r_j({\bf x})+S+\ve^{1.6}}\right]
    &\ll \E_{{\bf x}}
	\left[ \frac{\ve^{2.65}+r_j({\bf x})\cdot \sigma\cdot\sqrt{\sum_{i=1}^k \frac{1}{\al_i^2}}}{r_j({\bf x})}\cdot \frac{\min(r_j({\bf x}),S)}
	{r_j({\bf x})+S+\ve^{1.6}}\right]\\
    &\ll  \eps^{1.05}+\sigma\Expect{{\bf x}}{\sqrt{\sum_{i=1}^k \frac{1}{\bm{\al}_i^2}}\cdot \frac{\min(r_j({\bf x}),S)}{r_j({\bf x})+S+\ve^{1.6}}},
    \end{align*}
    and it is enough to bound the second term. Note that while we expect that each $\bm{\alpha}_i$ to be of the order $\log n/n$, convexity works against
    us and it could still be the case that $\sum\limits_{i=1}^{k}\frac{1}{\bm{\alpha}_i^2}$ could be large.
    The point is that in this case, some $\bm{\alpha}_i$ must be close to $0$, in which case $g_j({\bf x}_i)$ is very small
    -- cubically with $\bm{\alpha}_i$ -- thereby balancing the $1/\bm{\alpha}_i^2$ term. The following proposition formalizes this intuition,
    and the proof is deferred to Section~\ref{sec:def_proofs}
    \begin{proposition}\label{prop:convex_error}
      There is an absolute constant $A>0$ such that for any $z>0$ and $r\leq 1$ such that $r_j(x) = r\cdot g_j(x_i)$,
      it holds that
      \begin{equation*}
	\Expect{{\bf x}_i}{\sqrt{z+\frac{1}{\bm{\al}_i^2}} \cdot\frac{\min(r\cdot g_j({\bf x}_i),S)}
	{r\cdot g_j({\bf x}_i)+S+\ve^{1.6}}}
    \leq
	\Expect{{\bf x}_i}{\sqrt{z+A\frac{n^2}{\log^2 n}} \cdot \frac{\min(r\cdot g_j({\bf x}_i),S)}{r\cdot g_j({\bf x}_i)+S+\ve^{1.6}} +A\frac{n}{\log n} \cdot
\mathbbm{1}_{r\cdot g_j({\bf x}_i)\ge S}}.
	\end{equation*}
    \end{proposition}

    Applying Proposition~\ref{prop:convex_error} iteratively $k$ times (once for each $i$, taking $r = \prod\limits_{i'\neq i}{g_{j}(x_{i'})}$
    and the appropriate $z$), we get that
    \[
    \Expect{{\bf x}}{\sqrt{\sum_{i=1}^k \frac{1}{\bm{\al}_i^2}}\cdot \frac{\min(r_j({\bf x}),S)}{r_j({\bf x})+S+\ve^{1.6}}}
    \leq \Expect{{\bf x}}{\sqrt{k\cdot A\frac{n^2}{\log^2 n}}\cdot \frac{\min(r_j({\bf x}),S)}{r_j({\bf x})+S+\ve^{1.6}}+k\cdot A\frac{n}{\log n} \cdot
    \mathbbm{1}_{r_j({\bf x})\geq S}}.
    \]
    The proof is concluded by noting that $\frac{\min(r_j(x),S)}{r_j(x)+S+\ve^{1.6}}\leq \frac{r_j(x)}{r_j(x)+S}$.
    \end{proof}


We are now ready to finish the proof of inequality~\eqref{eq:mainA1}.
\subsubsection*{Proof of inequality~\eqref{eq:mainA1}}
Let $E$ be the event that:
(1) the support of $p({\bf x})$ has size at least $2$,
(2) $\sum\limits_{j} r_j({\bf x})\geq \eps^{1.6}$ and also for ${\bf x}+\bm{\De}$,
and (3) $\card{\bm{\De}_i}\leq \eps^{0.95}$ for all $i\in[n]$.
As we argued in Claims~\ref{cl:A-1},~\ref{cl:A-3},~\ref{cl:A-2} the contribution $({\bf x},\bm{\De})\not\in E$ to
the LHS of inequality~\eqref{eq:mainA1} is $\ll \eps$, hence it is enough to analyze the contribution of $({\bf x},\bm{\De})\in E$.


Denote $T(x) = \sum\limits_{j\in [m]} r_j(x)$.
\begin{align*}
&\Expect{{\bf x},\bm{\De}}{\norm{p({\bf x}) - p({\bf x}+\bm{\De})}_1 1_E}
= \Expect{{\bf x},\bm{\De}}{\sum\limits_{j\in [m]}\card{\frac{r_j({\bf x})}{T({\bf x})}-\frac{r_j({\bf x}+\bm{\De})}{T({\bf x}+\bm{\De})}} 1_E}\\
&= \underbrace{\Expect{{\bf x},\bm{\De}}{\sum\limits_{j\in [m]}\card{\frac{r_j({\bf x})}{T({\bf x})}-\frac{r_j({\bf x}+\bm{\De})}{T({\bf x}+\bm{\De})}}
1_E 1_{r_j({\bf x}) \leq \eps^{2.7}}}}_{(\rom{1})}
+\underbrace{\Expect{{\bf x},\bm{\De}}{\sum\limits_{j\in [m]}\card{\frac{r_j({\bf x})}{T({\bf x})}-\frac{r_j({\bf x}+\bm{\De})}{T({\bf x}+\bm{\De})}}
1_E 1_{r_j({\bf x}) >\eps^{2.7}}}}_{(\rom{2})}
\end{align*}
First, we show that $(\rom{1})\ll \eps$. As $T({\bf x})\geq \eps^{1.6}$ (since $E$ holds) and $r_j({\bf x})\leq \eps^{2.7}$, we get that
$r_j({\bf x})/T({\bf x})\leq \eps^{1.1}$, and next we argue that $r_j({\bf x}+\bm{\De})/T({\bf x}+\bm{\De})\ll \eps^{1.05}$.
Fix $j$ and suppose ${\bf x}_1,\ldots,{\bf x}_{k_j}$ are the ${\bf x}_i$'s that fall inside $I_j$. The following easy fact will be helpful.

\begin{fact}\label{fact:general}
  For all $x,\De$ we have
  $r_j(x+\De) = \sum\limits_{\substack{S\subseteq[k_j]}}{
  \prod\limits_{r\in S}{g_j(x_r)}
  \prod\limits_{r\not\in S}{(g_j(x_r)-g_j(x_r+\De_r))}}$.
\end{fact}
\begin{proof}
  Write $r_j(x+\De) =
  \prod\limits_{r=1}^{k_j} g_j(x_r+\De_r)
  =\prod\limits_{r=1}^{k_j}\left(g_j(x_r) + (g_j(x_r+\De_r)-g_j(x_r))\right)$
  and expand out.
\end{proof}

Combining Fact~\ref{fact:general} and Fact~\ref{fact:first_order_g}, we get that
\begin{align*}
r_j(x+\De)
\leq \sum\limits_{S\subseteq[k_j]}{
  \prod\limits_{r\in S}{g_j(x_r)}
  \prod\limits_{r\not\in S}{\card{g_j(x_r)-g_j(x_r+\De_r)}}}
&\leq \sum\limits_{S\subseteq[k_j]}{
\prod\limits_{r\in S}{g_j(x_r)}
B^{\card{S}} n^{3\card{S}}\prod\limits_{r\not\in S}(\alpha_r^3 + \card{\De_r}^3)}\\
&\leq \sum\limits_{S\subseteq[k_j]}{
\prod\limits_{r\in S}{g_j(x_r)}
B^{\card{S}} n^{3\card{S}}\prod\limits_{r\not\in S}\alpha_r^3}\\
&\phantom{\leq_S}+4^{n}B^{\card{n}}n^{3n}\max_{r}\card{\De_r}^3.
\end{align*}
Consider the right hand side above. For the first term we use $\alpha_r^3\ll g_j(x_r)$ to get it is at most
\[
\sum\limits_{S\subseteq[k_j]}{B'^{\card{S}} n^{3\card{S}} r_j(x_r)}
\leq (B'')^n n^{3n} \eps^{2.7} \leq \eps^{2.65}/2.
\]
For the second term we use $\card{\De_r}\leq \eps^{0.95}$ to bound it by $\eps^{2.65}/2$ as well.
We thus get $r_j(x+\De)\leq \eps^{2.65}$, and so $r_j(x+\De)/T(x+\De)\leq \eps^{1.05}$. Combined, we get that
\[
{(\rom{1})}\leq m(\eps^{1.1}+ \eps^{1.05})\ll \eps.
\]

Next, we handle $(\rom{2})$. Denote $T'(x) = \sum\limits_{j}{r_j(x) 1_{r_j(x)\geq \eps^{2.7}}}$, and note that
$T'(x)\geq T(x) - m\eps^{2.7}\geq (1-m\eps^{1.1})T(x)$ and similarly for $T'(x+\De)$. Thus, we may replace
$T(x), T(x+\De)$ with $T'(x), T'(x+\De)$ and incur (by the triangle inequality) a loss of at most $m\eps^{1.1}\ll \eps$.
Thus, we want to upper bound
\[
\underbrace{\Expect{{\bf x},\bm{\De}}{\sum\limits_{j\in [m]}\card{\frac{r_j({\bf x})}{T'({\bf x})}-\frac{r_j({\bf x}+\bm{\De})}{T'({\bf x}+\bm{\De})}} 1_E 1_{r_j({\bf x}) >\eps^{2.7}}}}_{(\rom{3})}\ll \eps.
\]
We intend to apply Claim~\ref{cl:A-6} with $r_j = r_j(x)$ and $d_j = r_j(x+\De) - r_j(x)$ for each $x$ separately, but for that
we first have to argue that $\card{d_j}\leq r_j/2$.
For each $i\in [n]$ there is $j$ such that $x_i\in I_j$, and we denote
$\al_i =\text{dist}\left(x_i, [z_j-\frac{\log n}{50n},z_j+\frac{\log n}{50n}]\right)$. Note that
\[
\eps^{2.7}\leq r_j(x) \leq g_j(x_i)\ll \left(\frac{n}{\log n}\alpha_i\right)^3,
\]
hence $\alpha_i\gg \frac{\log n}{n}\eps^{0.9}$, and for small enough $\eps$ we get that $\alpha_i\geq \eps^{0.91}\geq 2\card{\De_i}$. Therefore,
Combining Fact~\ref{fact:general} and Fact~\ref{fact:first_order_g} we get
\[
\card{d_j(x)}
=\card{r_j(x) - r_j(x+\De)}
\leq \sum\limits_{\substack{S\subseteq[k_j]\\S\neq[k_j]}}{
  \prod\limits_{r\in S}{g_j(x_r)}
  \prod\limits_{r\not\in S}{\card{g_j(x_r)-g_j(x_r+\De_r)}}}
\leq
\sum\limits_{\substack{S\subseteq[k_j]\\S\neq[k_j]}}{ B^{\card{S}} r_j(x)
\prod\limits_{r\not\in S}{\frac{\card{\De_i}}{\alpha_i}}}.
\]
Bounding $\frac{\card{\De_i}}{\alpha_i}\leq \eps^{0.95}/\eps^{0.91} = \eps^{0.04}$ we get that
\[
\card{r_j(x) - r_j(x+\De)}\leq r_j(x) \eps^{0.04}
\sum\limits_{\substack{S\subseteq[k_j]\\S\neq[k_j]}}{B^{\card{S}}}
\leq B'^n \eps^{0.04} r_j(x)
\leq r_j(x)/2 = r_j/2.
\]

Therefore, we may apply Claim~\ref{cl:A-6} and get that
\begin{align}
(\rom{3})&\ll \Expect{{\bf x},\bm{\De}}{\sum\limits_{j=1}^{m} \frac{|r_j({\bf x}) - r_j({\bf x}+\bm{\De})|}{r_j({\bf x})} \cdot \frac{\min(r_j({\bf x}), T'({\bf x}))}{T'({\bf x})}
1_E 1_{r_j({\bf x}) >\eps^{2.7}}} \notag\\\label{eq:2}
&\ll \Expect{{\bf x},\bm{\De}}{\sum\limits_{j=1}^{m} \frac{|r_j({\bf x}) - r_j({\bf x}+\bm{\De})|}{r_j({\bf x})} \cdot \frac{\min(r_j({\bf x}), T'({\bf x}))}{T'({\bf x})+\eps^{1.6}}1_E 1_{r_j({\bf x}) >\eps^{2.7}}},
\end{align}
where the last inequality holds since $T'({\bf x})\gg \eps^{1.6}$. Next, we wish to discard ${\bf x}_i$ that are very far from their closest center $z_j$.
For each $j$, note that $\left[z_j-\frac{\log n}{50 n},z_j+\frac{\log n}{50 n}\right]$ is exactly the set of $y$'s on which
$g_j(y) = 0$, and let $R_j\subseteq I_j$ be $R_j = \left[z_j-\frac{\log n}{25 n}-\eps^{0.95},z_j+\frac{\log n}{25n}+\eps^{0.95}\right] \setminus
\left[z_j-\frac{\log n}{50 n},z_j+\frac{\log n}{50 n}\right]$.
Note that for each $y\in I_j\setminus R_j$, we have that either $g_j(y) = 0$ if $y\in \left[z_j-\frac{\log n}{50 n},z_j+\frac{\log n}{50 n}\right]$, and otherwise $g_j(y) = 1$.
Furthermore, in the latter case we also have that $g_j(y+\bm{\De}_i) = 1$ since $\card{\bm{\De}_i}\leq \eps^{0.95}$.

We sample ${\bf x}$ in the following way. First, sample ${\bf t}_1,\ldots,{\bf t}_m$ the number of ${\bf x}_i$'s in each interval $I_1,\ldots,I_m$,
then for each $j$ sample ${\bf k}_j$ to be the number of $x_i$'s inside the interval $I_j$ that fall inside $R_j$.
Finally, for each $j\in[m]$ sample ${\bf k}_j$ points uniformly from $R_j$, ${\bf t}_j - {\bf k}_j$ uniformly from $I_j\setminus R_j$, and let ${\bf x}$
be the (multi-)set of all the sampled points. We condition on the ${\bf t}_j$'s and ${\bf k}_j$'s henceforth in~\eqref{eq:2}. Furthermore, we condition
on the identity of the $i$'s for which ${\bf x}_i\in I_j$ for each $j$.

Since $i$'s for which ${\bf x}_i\in I_j\in [z_j-\frac{\log n}{25 n}-\eps^{0.95},z_j+\frac{\log n}{25n}+\eps^{0.95}]$
do not affect both $r_j({\bf x})$ and $r_j({\bf x}+\bm{\De})$, we may ignore them and hence take expectation only over $i$'s such
${\bf x}_i\in R_j$. Call these $y$'s. Then from~\eqref{eq:2} we get
\begin{align*}
(\rom{3})
&\ll
\Expect{\vec{{\bf t}},\vec{{\bf k}}}
{
\Expect{{\bf y},\bm{\De}}
{
\sum\limits_{j=1}^{m} \frac{|r_j({\bf y}) - r_j({\bf y}+\bm{\De})|}{r_j({\bf y})} \cdot \frac{\min(r_j({\bf y}), T'({\bf y}))}{T'({\bf y})+\eps^{1.6}}
}}\\
&\leq
\Expect{\vec{{\bf t}},\vec{{\bf k}}}
{
\sum\limits_{j=1}^{m}
\Expect{{\bf y},\bm{\De}}
{\frac{|r_j({\bf y}) - r_j({\bf y}+\bm{\De})|}{r_j({\bf y})} \cdot \frac{\min(r_j({\bf y}), T'_{-j}({\bf y}))}{r_j({\bf y})+T'_{-j}({\bf y})+\eps^{1.6}}
}},
\end{align*}
where $T'_{-j}(x) = \sum\limits_{j'\neq j}r_{j'}(x) 1_{r_{j'}(x)\geq \eps^{2.7}}$.
Note that conditioned on $\vec{{\bf t}}= \vec{t},\vec{{\bf k}} = \vec{k}$, the values of ${\bf y}_i$'s such that ${\bf y}_i\in I_j$ are independent of $T'_{-j}({\bf y})$,
and they are distributed uniformly over $R_j$. Therefore, using Claim~\ref{cl:A-5} we have
\begin{align*}
(\rom{3})
&\ll
\Expect{\vec{{\bf t}},\vec{{\bf k}}}{
\sum\limits_{j=1}^{m}
\eps^{1.05}+ {\bf k}_j\frac{\sigma n}{\log n}  \cdot \cProb{{\bf y}}{\vec{{\bf t}},\vec{{\bf k}}}{r_j({\bf y})\ge T'_{-j}({\bf y})}
+\Expect{{\bf y}}{\sigma \frac{n}{\log n}\sqrt{{\bf k}_j}\frac{r_j({\bf y}) }{T'({\bf y})}}
}.\\
&\leq
m\eps^{1.05}
+ n^2\sigma
\sum\limits_{j=1}^{m}\Prob{y}{r_j(y)\ge T'_{-j}(y)}
+ \sigma\frac{n}{\log n}
\Expect{\vec{t},\vec{k}}{\sqrt{\max_{j} k_j}}.
\end{align*}
Note that if $T'_{-j}(x)\leq r_j(x)$, then
\[
T(x) \leq
T'_{-j}(x) + r_j(x) + \sum\limits_{j'}{r_{j'}(x) 1_{r_{j'}(x)\leq \eps^{2.7}}}
\leq 2r_j(x) + m\cdot \eps^{2.7}
\leq 3,
\]
so we bound the sum on the right hand side by $m\Prob{x}{T(x)\leq 3}$. For the expectation,
we use Cauchy-Schwarz and overall we get
\[
(\rom{3})
\leq m\eps^{1.05} + n^3\sigma \Prob{{\bf x}}{T({\bf x})\leq 3} +
\sigma\frac{n}{\log n}\sqrt{\Expect{\vec{{\bf t}},\vec{{\bf k}}}{\max_{j} {{\bf k}_j}}}.
\]
The first term is clearly $\ll \eps$. For the second term, we use Claim~\ref{cl:A-7} below, that asserts that
$\Prob{{\bf x}}{T({\bf x})\leq 3}\leq n^{-\omega(1)}$, hence by the definition of $\sigma$ the second term is
also $\ll \eps$. For the third term, note that each ${\bf k}_j$ is a sum of $n$ independent Berounlli random variables
with parameter $p\leq\log n/n$, therefore by Chernoff bound
\[
\Prob{}{{\bf k}_j\geq 10\log n}
\leq e^{-\frac{1}{3}9^2 \log n}
\leq n^{-9}.
\]
The union bound now implies that $\Prob{}{\max_j {\bf k}_j\geq 10\log n}\leq n^{-8}$, and hence
\[
\Expect{\vec{{\bf t}},\vec{{\bf k}}}{\max_{j} {\bf k}_j}
\leq n^{-8}\cdot n + 10\log n\ll \log n.
\]
Using the definition of $\sigma$, we get that the third term is also $\ll \eps$. Combining all,
we get that $(\rom{3})\ll \eps$, and we are done.
\qed

\begin{claim}\label{cl:A-7}
\begin{equation}
\Pr_{{\bf x}} \left[\sum_j r_j({\bf x}) \le 3\right] < n^{-\omega(1)}.
\end{equation}
\end{claim}

\begin{proof}
	The proof is very similar to the analysis of Case~(B) above. In particular, similarly to inequality~\eqref{eq:notEi},
	$$
	\Pr[r_j({\bf x})<1] = \left( 1-\frac{2m \log n}{25 n} \right)^{t_j}
    \geq
	\left( 1-\frac{2m \log n}{25 n} \right)^{2 \cdot n/m} > e^{-2 \log n/25} = n^{-2/25},
	$$
	as long as $t_j < 2\cdot n/m$ (which is the case except with probability $n^{-\omega(1)}$. Since $m > n^{2/25} \cdot n^{\Omega(1)}$, the probability of not having at least three $r_j(x)$'s equal to $1$ is $n^{-\omega(1)}$.
\end{proof}















































































































































































































































































































































































































































































































































































































































































































































































































































































































































































\setcounter{thm}{57}\appendix
\section{Deferred proofs}\label{sec:def_proofs}
\subsection{Proof of Claim~\ref{cl:A-6}}
We split the proof into two cases.
	\paragraph{Case 1: $r_i \le T/2$ for all $i$.} In this case,
	$\min(r_i, T-r_i)=r_i$ for all $i$, and the sum on the RHS of \eqref{eq:cl:A-6:0} is just $(\sum_i |d_i|)/T$. We have
	\begin{align*}
	\norm{p-q}_1 = \sum_i \left| \frac{r_i+d_i}{T'} - \frac{r_i}{T} \right|
    &\le
	\sum_i \left| \frac{r_i+d_i}{T} - \frac{r_i}{T} \right| +
	\sum_i \left| \frac{r_i+d_i}{T'} - \frac{r_i+d_i}{T} \right| \\
	&=\sum_i \frac{|d_i|}{T} + \left|\frac{1}{T'}-\frac{1}{T}\right|\cdot
	\sum_i (r_i+d_i)\\
    &= \sum_i \frac{|d_i|}{T}  + \left|1-\frac{T'}{T}\right|\\
    &=\sum_i \frac{|d_i|}{T} +\frac{1}{T} \cdot \left| \sum_i d_i \right| \\
	&\leq 2\cdot \sum_i \frac{|d_i|}{T}.
	\end{align*}
	
	\paragraph{Case 2: one of the $r_i$'s is greater than $T/2$.}
	Without loss of generality, $r_1>T/2$. Denote by $S:=\sum_{i>1} r_i = T-r_1$; $S':=\sum_{i>1} (r_i+d_i) = T'-r_1-d_1$.
	 In this case, the RHS of \eqref{eq:cl:A-6:0} is given by
	 \begin{equation}
	 \label{eq:cl:A-6:1}
	 \frac{|d_1| \cdot S}{r_1\cdot T} + \sum_{i>1} \frac{|d_i|}{T}.
	 \end{equation}
	 We will estimate $|p_1-q_1|$ and $\sum_{i>1} |p_i-q_i|$ separately. First,
     note that $T'\geq T - \sum_{j}\card{d_j}\geq T/2$.

For $|p_1-q_1|$, we have	
	 \begin{align*}
	 |p_1-q_1|= \left| \frac{r_1}{T} - \frac{r_1+d_1}{T'} \right| =
	  \left| \frac{r_1\cdot(S'-S)+ d_1\cdot S}{T\cdot T' } \right|
\le
	   2\left| \frac{S'-S}{T} \right|+
	   2\left| \frac{ d_1\cdot S}{T\cdot r_1 } \right|
\le \sum_{i>1} \frac{|d_i|}{T} +
	    \frac{|d_1| \cdot S}{r_1\cdot T}.
	 \end{align*}
In the third transition, we used the fact that $T' \geq T/2\geq r_1/2$.
	
	 For $\sum_{i>1} |p_i-q_i|$, by a similar calculation to the first case we have
	 \begin{align*}
	 \sum_{i>1}|p_i-q_i| =  \sum_{i>1} \left| \frac{r_i+d_i}{T'} - \frac{r_i}{T} \right|
     &\le
	 \sum_{i>1} \left| \frac{r_i+d_i}{T} - \frac{r_i}{T} \right| +
	 \sum_{i>1} \left| \frac{r_i+d_i}{T'} - \frac{r_i+d_i}{T} \right|\\
     &\leq
	 \sum_{i>1} \frac{|d_i|}{T}+\left|\frac{1}{T'}-\frac{1}{T}\right|\cdot
	 \sum_{i>1} r_i+d_i
     = \sum_{i>1} \frac{|d_i|}{T} + \card{\frac{1}{T'}-\frac{1}{T}} S',
     \end{align*}
     and it is enough to bound $\card{\frac{1}{T'}-\frac{1}{T}} S'$ by constant times the expression in~\eqref{eq:cl:A-6:1}.
     We have

  \begin{align*}
     \card{\frac{1}{T'}-\frac{1}{T}} S'=
	 \left|\frac{S'\cdot (S'-S) + S'\cdot d_1}{T'T}\right|
     &\leq
	 \left|\frac{(S'+d_1)\cdot (S'-S)}{T'T}\right|+
	 \left|\frac{ S\cdot d_1}{T'T}\right|\\
     &\leq
     \left|\frac{S'\cdot (S'-S)}{T'T}\right|
     +\left|\frac{d_1\cdot (S'-S)}{T'T}\right|
	 + 2\cdot
	 \left|\frac{ S\cdot d_1}{T^2}\right|,
	 \end{align*}
where in the last transition we used $T' \geq T/2 > 0$.
We bound each term separately.
For the first term, as $T'\geq T/2$, $\card{S'}\leq 2T$ (since $\card{d_i}\leq r_i$) we get
\[
\card{\frac{S'\cdot (S'-S)}{T'T}}
\leq 4\card{\frac{S'-S}{T}}
\leq 4\sum\limits_{i\geq 2}{\frac{\card{d_i}}{T}}.
\]

For the second term, we have $\card{d_1}\leq r_1\leq T$, $T'\geq T/2$ and so
\[
\card{\frac{d_1\cdot (S'-S)}{T'T}}
\leq 2\frac{\card{S'-S}}{T}
\leq 2\sum\limits_{i\geq 2}{\frac{\card{d_i}}{T}}.
\]

For the third term, we have, as $T\geq r_1$,
$\frac{ S\cdot d_1}{T^2}
\leq \frac{\card{d_1}}{r_1}\frac{S}{T}$.

\subsection{Proof of Proposition~\ref{prop:convex_error}}
	We will use the fact for points $x_i$ in our domain, $g_j(x_i) \ee \left(\frac{n}{\log n}\alpha_i\right)^3$.
	We consider two cases, based on the values of $S$ and $r$.
	\paragraph{Case 1: $\Pr_{{\bf x}_i} [r\cdot g_j({\bf x}_i)>S] < 1/2$.}
    We claim that for a sufficiently large constant $A>0$,
		\begin{equation*}
	\underbrace{\Expect{{\bf x}_i}{\sqrt{z+\frac{1}{\bm{\al}_i^2}} \cdot\frac{\min(r\cdot g_j({\bf x}_i),S)}
	{r\cdot g_j({\bf x}_i)+S+\ve^{1.6}}}}_{(\rom{1})}
    \leq
	\underbrace{\E_{{\bf x}_i}\left[ \sqrt{z+An^2/\log^2 n} \cdot \frac{\min(r\cdot g_j({\bf x}_i),S)}
	{r\cdot g_j({\bf x}_i)+S+\ve^{1.6}}  \right]}_{(\rom{2})}
    .
	\end{equation*}
	To do that, we compare both sides to $\Expect{{\bf x}_i}{\sqrt{z} \cdot\frac{\min(r\cdot g_j({\bf x}_i),S)}{r\cdot g_j({\bf x}_i)+S+\ve^{1.6}}} $.	
	For $(\rom{1})$, we have
    \[
	\Expect{{\bf x}_i}{\left(\sqrt{z+1/\bm{\al}_i^2} -\sqrt{z}\right)\cdot \frac{\min(r\cdot g_j({\bf x}_i),S)}
	{r\cdot g_j({\bf x}_i)+S+\ve^{1.6}}}
    \ll \Expect{{\bf x}_i}{\frac{1/\bm{\alpha}_i^2}{\sqrt{z+1/\bm{\al}_i^2}}\cdot \frac{\min(r\cdot g_j({\bf x}_i),S)}
	{r\cdot g_j({\bf x}_i)+S+\ve^{1.6}}}.
	\]
    Since $\bm{\alpha}_i\ll \log n/n$ always, we may further upper bound this by
    \[
    \ll \Expect{{\bf x}_i}{\frac{1/\bm{\alpha}_i^2}{\sqrt{z+A/(\log n/n)^2}}\cdot \frac{\min(r\cdot g_j({\bf x}_i),S)}
	{r\cdot g_j({\bf x}_i)+S+\ve^{1.6}}}
    \ll \Expect{{\bf x}_i}{\frac{1/\bm{\alpha}_i^2}{\sqrt{z+An^2/\log^2 n}}\cdot \frac{r\left(\frac{n}{\log n}\bm{\al}_i\right)^3}
	{S+\ve^{1.6}}},
	\]
    where we used $\min(r\cdot g_j({\bf x}_i),S)\leq rg_j({\bf x}_i)$ and the asymptotic we have for $g_j$. Simplifying and using
    $\Expect{{\bf x}_i}{\bm{\alpha}_i}\ll \log n/n$, we get that the last expression is equal to
    \[
    \frac{n^2}{\log^2 n}\frac{1}{\sqrt{z+An^2/\log^2 n}}\cdot \frac{r}{S+\ve^{1.6}}.
    \]


    For $(\rom{2})$, we have
    \[
    \Expect{{\bf x}_i}{\left(\sqrt{z+An^2/\log^2 n}-\sqrt{z}\right)\cdot \frac{\min(r\cdot g_j({\bf x}_i),S)}{r\cdot g_j({\bf x}_i)+S+\ve^{1.6}}}
    \gg
     \Expect{{\bf x}_i}{\frac{An^2/\log^2 n}{\sqrt{z+An^2/\log^2 n}}\cdot \frac{\min(r\cdot g_j({\bf x}_i),S)}{r\cdot g_j({\bf x}_i)+S+\ve^{1.6}}}.
    \]
    Restricting to the event $E$ that $r g_j({\bf x}_i)\leq S$ (that has probability at least $1/2$ by assumption), we have that the last expression
    is at least
    \[
    \gg\cExpect{{\bf x}_i}{E}{\frac{An^2/\log^2 n}{\sqrt{z+An^2/\log^2 n}}\cdot \frac{r\cdot g_j({\bf x}_i)}{S+\ve^{1.6}}}
    \gg \frac{An^2/\log^2 n}{\sqrt{z+An^2/\log^2 n}}\cdot \frac{r}{S+\ve^{1.6}},
    \]
    where the last inequality holds since $\cExpect{\alpha_i}{E}{g_j(x_i)}\gg 1$ (this is true for any event $E$ with constant probability
    in our range of interest of ${\bf x}_i$'s). Combining the bounds for $(\rom{1}),(\rom{2})$, we see that we may pick large enough $A$ so
    that
    \[
    \Expect{{\bf x}_i}{\left(\sqrt{z+1/\al_i^2} -\sqrt{z}\right)\cdot \frac{\min(r\cdot g_j({\bf x}_i),S)}
	{r\cdot g_j({\bf x}_i)+S+\ve^{1.6}}}
\leq \Expect{{\bf x}_i}{\left(\sqrt{z+An^2/\log^2 n}-\sqrt{z}\right)\cdot \frac{\min(r\cdot g_j({\bf x}_i),S)}{r\cdot g_j({\bf x}_i)+S+\ve^{1.6}}},
    \]
    and hence $(\rom{1})\leq (\rom{2})$. Let $A_1$ be a large enough value of $A$ so that this holds.

    \paragraph{Case 2: $\Prob{{\bf x}_i}{r\cdot g_j({\bf x}_i)>S} \geq 1/2$.}
    Using $\sqrt{a+b}\leq \sqrt{a}+\sqrt{b}$, we have
    \[
    \Expect{{\bf x}_i}{\sqrt{z+\frac{1}{\bm{\al}_i^2}} \cdot\frac{\min(r\cdot g_j({\bf x}_i),S)}
	{r\cdot g_j({\bf x}_i)+S+\ve^{1.6}}}
    \leq
    \underbrace{\Expect{{\bf x}_i}{\sqrt{z} \cdot\frac{\min(r\cdot g_j({\bf x}_i),S)}
	{r\cdot g_j({\bf x}_i)+S+\ve^{1.6}}}}_{(\rom{3})}
    +\underbrace{\Expect{{\bf x}_i}{\frac{1}{\bm{\alpha}_i} \cdot\frac{\min(r\cdot g_j({\bf x}_i),S)}
	{r\cdot g_j({\bf x}_i)+S+\ve^{1.6}}}}_{(\rom{4})}.
    \]
    Clearly,
    $(\rom{3})\leq \Expect{{\bf x}_i}{\sqrt{z+A\frac{n^2}{\log^2 n}} \cdot \frac{\min(r\cdot g_j({\bf x}_i),S)}{r\cdot g_j({\bf x}_i)+S+\ve^{1.6}}}$,
    and we upper bound $(\rom{4})$. Recall that $g_j({\bf x}_i) \ee \left(\frac{n}{\log n} \bm{\alpha}_i\right)^3$, so
    \[
    (\rom{4})
    \ll \Expect{{\bf x}_i}{\frac{1}{\bm{\alpha}_i} \cdot
    \frac{\min(r(n\bm{\alpha}_i/\log n)^3,S)}{B\cdot r (n\bm{\al}_i/\log n)^3+S+\ve^{1.6}}},
    \]
    for some absolute constant $B>0$. Writing the last expression as an integral, we note
    that $\bm{\alpha}_i$ is distributed uniformly on the interval $[0,\frac{\log n}{50n}+\eps^{0.95}]$, so
    we get
    \[
    (\rom{4})\ll\left(\frac{n}{\log n}\right)\int_{0}^{\frac{\log n}{25 n}}\frac{1}{t}\frac{\min(r(n t/\log n)^3,S)}{B\cdot r (nt/\log n)^3+S+\ve^{1.6}} dt.
    \]
    We break the range of integration into $R_1 = \left[0, (S/r)^{1/3} \frac{\log n}{n}\right]$,
    and $R_2 = \left[(S/r)^{1/3} \frac{\log n}{n},\frac{\log n}{25 n}\right]$. On $R_1$ our expression is equal to
    \[
    \left(\frac{n}{\log n}\right)^2\int_{0}^{\left(\frac{S}{r}\right)^{1/3} \frac{\log n}{n}}\frac{r(n t/\log n)^2}{B\cdot r (nt/\log n)^3+S+\ve^{1.6}} dt
    \ll
    \left(\frac{n}{\log n}\right)^4\int_{0}^{\left(\frac{S}{r}\right)^{1/3} \frac{\log n}{n}}\frac{r t^2}{S} dt
    \ll \frac{n}{\log n}.
    \]
%
%
%
%
%
%
%

    On $R_2$ our expression is at most
    \[
    \left(\frac{n}{\log n}\right)\int_{\left(\frac{S}{r}\right)^{1/3} \frac{\log n}{n}}^{\frac{\log n}{25n}}
    \frac{1}{t}\frac{S}{B\cdot r (nt/\log n)^3} dt
    \ll\frac{S}{r}\left(\frac{\log n}{n}\right)^2\int_{\left(\frac{S}{r}\right)^{1/3} \frac{\log n}{n}}^{\frac{\log n}{25n}}
    \frac{1}{t^4} dt.
    \]
    Computing the integral, we see it is at most $(\left(\frac{S}{r}\right)^{1/3} \frac{\log n}{n})^{-3}$, hence
    the overall expression is $\ll n/\log n$, and since $\Expect{}{\mathbbm{1}_{r\cdot g_j(x_i)>S}}\geq 1/2$ we
    conclude that there is $A_2>0$ such that
    \[
    (\rom{4})\leq A_2 \frac{n}{\log n} \Expect{{\bf x}_i}{\mathbbm{1}_{r\cdot g_j({\bf x}_i)>S}}.
    \]

    The proposition is thus proven for $A = \max(A_1,A_2)$\qed

    \section{From Noise Sensitivity to Surface Area}\label{sec:ns_sa_reduce}
    Let $D_{\vec{r}}$ be a family of tilings of $\mathbb{R}^n$ that are constructed from Lemma~\ref{lem:main_ub_use}.
    I.e., the family $D_{\vec{r}}$ satisfies that the there is $A = O(n/\sqrt{\log n})$ such that for sufficiently small $\eps$, we have that
    \[
      \Expect{\vec{r}}{\Prob{\substack{{\bf x}\in D_{\vec{r}}\\ \bm{\De}\sim N(0,\eps^2\cdot I_n)}}
      {{\bf x},{\bf x}+\bm{\De}\text{ fall in different cells of the tiling induced by }D_{\vec{r}}}}\leq A\eps.
    \]
    Let $k_0$ be the first integer such that this condition holds for any $0 < \eps\leq 2^{-k_0}$.
    Thus, defining for each $k\geq k_0$ the set
    \[
    G_k = \Big\{\vec{r} \qquad\Big|\qquad \Prob{\substack{{\bf x}\in D_{\vec{r}}\\ \bm{\De}\sim N(0,4^{-k}\cdot I_n)}}
    {{\bf x},{\bf x}+\bm{\De}\text{ lie in different cells of the tiling of }S_{\vec{r}}}\leq 2\cdot A 2^{-k}\Big\},
    \]
    we have by Markov's inequality that $\Prob{\vec{r}}{\vec{r}\in G_k}\geq \half$.
    \begin{claim}
      The sets $G_k$ are monotone decreasing, i.e.\ for each $k$, $G_{k+1}\subseteq G_k$.
    \end{claim}
    \begin{proof}
      Fix $\vec{r}\in G_{k+1}$. Let $\De\sim N(0,4^{-k-1}\cdot I_n)$, and note that $\De' = 2\cdot\De\sim N(0,4^{-k}\cdot I_n)$.
      Thus,
      \begin{align}
      \Prob{\substack{{\bf x}\in D_{\vec{r}}\\ \bm{\De'}\sim N(0,4^{-k}\cdot I_n)}}{{\bf x},{\bf x}+\bm{\De'}\text{ in different cells}}
      &\leq \Prob{\substack{{\bf x}\in D_{\vec{r}}\\ \bm{\De}\sim N(0,A 4^{-k-1}\cdot I_n)}}{{\bf x},{\bf x}+\bm{\De}\text{ in different cells}}\notag\\\label{eq:apx_1}
      &+\Prob{\substack{{\bf x}\in D_{\vec{r}}\\ \bm{\De}\sim N(0,4^{-k-1}\cdot I_n)}}{{\bf x}+\bm{\De},{\bf x}+2\bm{\De}\text{ in different cells}}.
      \end{align}
      First, we argue that the second probability on the right hand side is equal to the first one. To see that,
      denote $y = x+\De$ and observe that
      the points $y,y+\De$ lie in different cells of the tiling induced by $D_{\vec{r}}$ if and only if the points
      $y\pmod{D_{\vec{r}}}$, $y\pmod{D_{\vec{r}}} + \De$ lie in different cells. Additionally, note for any fixed $\De$,
      the distribution of ${\bf y}\pmod{D_{\vec{r}}}$ when we take ${\bf x}\in_R D_{\vec{r}}$, is uniform over $D_{\vec{r}}$.

      Therefore, the bound we get from~\eqref{eq:apx_1} is (using the fact that $\vec{r}\in G_{k+1}$)
      \[
      2\cdot \Prob{\substack{{\bf x}\in D_{\vec{r}}\\ \bm{\De}\sim N(0,4^{-k-1}\cdot I_n)}}{{\bf x},{\bf x}+\bm{\De}\text{ in different cells}}\leq
      2\cdot 2\cdot A 2^{-(k+1)}
      = 2\cdot A 2^{-k},
      \]
      and so $\vec{r}\in G_{k}$.
    \end{proof}

    \begin{claim}\label{claim:plenty_of_tilings}
      It holds that $\Prob{\vec{r}}{\vec{r}\in \bigcap_{k\geq k_0} G_k}\geq \half$, and in particular $\bigcap_{k\geq k_0} G_k$
      is not empty.
    \end{claim}
    \begin{proof}
      Define the sequence of functions $g_m(\vec{r}) = \mathrm{1}_{\vec{r}\in \bigcap_{k_0\leq k\leq m} G_k}$,
      and also $g = \mathrm{1}_{\vec{r}\in \bigcap_{k\geq k_0} G_k}$. Clearly, on each $\vec{r}$, the sequence $g_m(\vec{r})$
      is monotonically decreasing to $g(\vec{r})$, and in other words we have monotone pointwise convergence of
      the non-negative functions $g_m$ to $g$.
      Thus, by the monotone convergence theorem
      \[
      \Prob{\vec{r}}{\vec{r}\in \bigcap_{k\geq 0} G_k}
      = \Expect{\vec{r}}{g(\vec{r})}
      = \Expect{\vec{r}}{\lim_{k\rightarrow\infty} g_k(\vec{r})}
      = \lim_{k\rightarrow\infty} \Expect{\vec{r}}{g_k(\vec{r})}.
      \]
      By the previous claim, $g_m = \mathrm{1}_{G_m}$, hence $\Expect{\vec{r}}{g_m(\vec{r})}\geq \half$
      and in particular the limit above is at least $\half$.
    \end{proof}

    Pick $\vec{r}^{\star}\in \bigcap_{k\geq k_0} G_k$, $\eps = 2^{-k_0}$ and denote $D = D_{\vec{r}^{\star}}$ for the rest of the proof. Clearly $D$ induces a tiling
    of the space $\mathbb{R}^n$, and next we will show that the surface area of $D$ is $O(A) = O(n/\sqrt{\log n})$, as desired.

    Towards this end, we will use Lemma~\ref{lem:santalo} that tells us that the surface
    area of $D$ is a constant multiple of
    \[
    \frac{1}{\eps}\Expect{\substack{{\bf x}\in_R D\\ \bm{\De}\sim N(0,\eps^2 I_n)}}{\card{({\bf x},{\bf x}+\bm{\De})\cap \partial D}},
    \]
    and we first observe that $({\bf x},{\bf x}+\bm{\De})\cap \partial D$ is almost surely countable.
    \footnote{The diligent reader may note that here, we are only considering intersections of the surface with the
    open interval $(x,x+\De)$ as opposed to the closed interval. This does not make any difference,
    since the contribution of the endpoints is proportional to the measure of $\partial D$. Hence, if the measure
    of $\partial D$ is $0$ they endpoints contribute $0$ to that expectation, and if the measure of $\partial D$
    is positive, then the expectation is infinite either way.}

    \begin{claim}\label{claim:countable}
      Let $\eps>0$ and sample ${\bf x}\in_R D$, $\bm{\De}\sim N(0,\eps^2 I_n)$. Then with probability $1$,
      $({\bf x},{\bf x}+\bm{\De})\cap \partial D$ is finite or countable.
    \end{claim}
    \begin{proof}
      Recall that by Lemma~\ref{lem:main_ub_use}, $D$ is a countable union of semi-algebraic sets, say $B_1,B_2,\ldots$.
      Note that for each semi-algebraic set $B_i$, the probability that $({\bf x},{\bf x}+\bm{\De})\cap \partial B_i$ is infinite is $0$, hence
      by the union bound, with probability $1$ all of these sets are finite, in which case $({\bf x},{\bf x}+\bm{\De})\cap \partial D$ is finite
      or countable.
    \end{proof}


    For a parameter $h$, a point $x\in\mathbb{R}^n$ and a direction $\De$, we say a
    point $y\in (x,x+\De)$ is $h$-isolated if
    \begin{enumerate}
      \item It holds that $y\in \partial D$.
      \item The neighbourhood of radius $h$ around $y$ does not contain $x,x+\De$ or any point from $\partial D$ (besides $y$).
    \end{enumerate}
    Define the quantity $g_{m}(x,\De)$ to be the number of $2^{-m}\norm{\Delta}_2$-isolated points in the interval $[x,x+\De]$.

    \begin{claim}\label{claim:converge}
      $g_{m}(x,\De)$ is an increasing sequence in $m$, and for any $x,\De$ for which Claim~\ref{claim:countable} holds, we have
      \[
      \lim_{m\rightarrow \infty} g_{m}(x,\De) = \card{(x,x+\De)\cap \partial D}.
      \]
    \end{claim}
    \begin{proof}
      The monotonicity of $g_m(x,\De)$ in $m$ is clear, and also that
      $g_{m}(x,\De)\leq \card{(x,x+\De)\cap \partial D}$. We set
      $\ell = g_{m}(x,\De)$ and split the rest of the proof according to
      whether $\ell$ is finite or not.

      \paragraph{Case 1: $\ell$ is finite.}
      In this case we argue that $g_{m}(x,\De) = \card{(x,x+\De)\cap \partial D}$
      for large enough $m$.  To see that, let $y_1,\ldots,y_\ell\in (x,x+\De)$ be all of the intersection
      points of $(x,x+\De)$ and $\partial D$, and take large enough $m$
      so that $2^{-m}\norm{\De}_2$ is smaller than all of the distances $\norm{y_i-y_j}_2$,
      $\norm{y_i-x}_2$, $\norm{y_i - (x + \De)}_2$ for all $i$ and $j$.

      \paragraph{Case 1: $\ell$ is infinite.}
      Consider the set $S = [x,x+\De]\cap \partial D$, and note that it is a closed.
      By Claim~\ref{claim:countable}, $S$ is countable, and we argue that $S$ must have an
      isolated point. Otherwise, $S$ is a closed set and has no isolated point, i.e.\
      it s a perfect set, but then it must be uncountable (e.g.\ see~\cite{perfect}).
      We thus conclude that $S$ has an isolated point $w_1$; we may remove it from
      $S$, have that the resulting set is again closed and countable, so we may again find
      an isolated point. Repeating this argument, for any $v\in\mathbb{N}$ we may find
      a collection of isolated points $w_1,\ldots,w_v\in S$ that are all different from
      $x$ and $x+\De$. As in case $1$, we conclude that $g_m(x,\De) \geq v$
      for large enough $m$, and since it holds for any $v$ we conclude that $\lim_{m\rightarrow \infty} g_m(x,\De) = \infty$.
    \end{proof}

    By Lemma~\ref{lem:santalo}, we have that the surface area of $D$ is at most a constant multiple of
    \[
    \frac{1}{\eps}\hspace{-0.5ex}\Expect{\substack{{\bf x}\in_R D\\ \bm{\De}\sim N(0,\eps^2 I_n)}}{\card{({\bf x},{\bf x}+\bm{\De})\cap \partial D}}
    =\frac{1}{\eps}\hspace{-0.5ex}\Expect{\substack{{\bf x}\in_R D\\ \bm{\De}\sim N(0,\eps^2 I_n)}}{\lim_{m\rightarrow \infty} g_{m}({\bf x},\bm{\De})}
    =\lim_{m\rightarrow \infty} \frac{1}{\eps}\hspace{-0.5ex}\Expect{\substack{{\bf x}\in_R D\\ \bm{\De}\sim N(0,\eps^2 I_n)}}{g_{m}({\bf x},\bm{\De})}.
    \]
    In the first transition we used Claims~\ref{claim:converge} and~\ref{claim:countable},
    and in the second one we used monotone convergence. Thus, if we assume that the surface area of
    $D$ is larger than $c\cdot A$ for a sufficiently large absolute constant $c$, then we get that
    $\lim_{m\rightarrow \infty} \frac{1}{\eps}\Expect{\substack{x\in_R D\\ \De\sim N(0,\eps^2 I_n)}}{g_{m}(x,\De)}\geq 10 A$.
    In the rest of the proof we will reach a contradiction and thereby show that for sufficiently large
    absolute constant $c$, the surface area of $D$ is at most $cA$, as required.

    By properties of limits, we conclude there exists $m$ such that
    \begin{equation}\label{eq:apx_2}
    \Expect{\substack{{\bf x}\in_R D\\ \bm{\De}\sim N(0,\eps^2 I_n)}}{g_{m}({\bf x},\bm{\De})}\geq 5 A\eps,
    \end{equation}
    and we fix this $m$ henceforth.

    Take $0<\delta\leq 2^{-m}$, and consider the following experiment. Take ${\bf x}\in_R D$ uniformly at random,
    $\bm{\De}\sim N(0,\eps^2 I_n)$ and take a uniformly random point ${\bf y}\in_R [{\bf x},{\bf x}+\bm{\De}]$. We
    consider the event $E$ in which the points ${\bf y}$ and ${\bf y}+\delta\bm{\De}$ lie in different cells in the tiling
    induced by $D$.

    \begin{claim}\label{claim:apx_one}
     For any $x,\De$ we have that $\cProb{{\bf y}}{x,\De}{E}\geq \delta g_m(x,\De)$.
    \end{claim}
    \begin{proof}
    Let $\ell = g_m(x,\De)$, and let $z_1,\ldots,z_{\ell}$ be the $2^{-m}\norm{\De}_2$-isolated points on the interval
    $(x,x+\De)$. For each $j$, let
    $I_j = (z_j - \delta\De,z_j)$, and note that as $\delta\leq 2^{-m}$ and the isolation of the points, we conclude that
    the intervals $I_j$ are disjoint and contained in $(x,x+\De)$.
    Also, note that if we pick $y\in I_j$, then $y$ and $y+\delta \De$ lie in different cells of the tiling induced by $D$;
    this holds since the interval between them contains exactly one point from $\partial D$ (namely, the point $z_j$). Therefore,
    \[
    \cProb{{\bf y}}{x,\De}{E}
    \geq\sum\limits_{j=1}^{\ell} \cProb{{\bf y}}{x,\De}{{\bf y}\in I_j}
    \geq\sum\limits_{j=1}^{\ell} \frac{\delta\norm{\De}_2}{\norm{\De}_2}
    =\delta\ell.\qedhere
    \]
    \end{proof}

   \begin{claim}\label{claim:apx_two}
     $\Prob{{\bf x},{\bf \De},{\bf y}}{E}\leq 2A\delta\eps$.
   \end{claim}
   \begin{proof}
      Consider ${\bf x},\bm{\De},{\bf y}$ the random variables in the definition of the event $E$. Let ${\bf z} = {\bf y}\pmod{D}$,
      and note that the points ${\bf y}$ and ${\bf y}+\delta\bm{\De}$ fall in different cells if and only if the points ${\bf z}$ and
      ${\bf z}+\delta\bm{\De}$ fall in different cells. Therefore, the probability of $E$ is exactly the probability that ${\bf z}$, ${\bf z}+\delta\bm{\De}$
      fall in different cells. Further, note that conditioned on $\bm{\De}$, the distribution of ${\bf z}$ is uniform over $D$, so
      \[
      \Prob{\bm{\De}\sim N(0,\eps^2 I_n)}{{\bf z}, {\bf z}+\delta\bm{\De}\text{ lie in different cells of } D}
      =\Prob{{\bf \De'}\sim N(0,\delta^2\eps^2 I_n)}{{\bf z}, {\bf z}+\bm{\De'}\text{ lie in different cells of }D},
      \]
      which is at most $2A\delta\eps$ by the choice of $D$ and the fact that $\delta\eps\leq \eps\leq 2^{-k_0}$.
    \end{proof}

    Combining the above claims we reach a contradiction:
    \[
    2A\delta\eps
    \geq
    \Prob{{\bf x},\bm{\De},{\bf y}}{E}
    =\Expect{{\bf x},\bm{\De}}{\cProb{{\bf y}}{{\bf x},\bm{\De}}{E}}
    \geq
    \Expect{{\bf x},\bm{\De}}{\delta g_m({\bf x},\bm{\De})}
    \geq \delta\cdot 5A\eps,
    \]
    and contradiction. The first transition is by Claim~\ref{claim:apx_one}, the second
    transition is by conditional probability formula, the third transition is by Claim~\ref{claim:apx_two}
    and the final one is by equation~\eqref{eq:apx_2}.\qed

















